\documentclass[11pt]{article}
\usepackage[margin=1.2in]{geometry}
\usepackage{amsmath,amssymb,amsthm}
\usepackage{xr}
\externaldocument{../paper2b_note/note}
\usepackage{hyperref}
\sloppy

\newtheorem{theorem}{Theorem}
\newtheorem{corollary}[theorem]{Corollary}
\newtheorem{lemma}[theorem]{Lemma}
\newtheorem{conjecture}[theorem]{Conjecture}
\newtheorem{proposition}[theorem]{Proposition}
\theoremstyle{remark}
\newtheorem{remark}[theorem]{Remark}

\newcommand{\E}{\mathbb{E}}
\newcommand{\Q}{\mathbb{Q}}

\title{The biregular case: supplementary material\\
to the weighted plane class and its obstructions}
\author{Vico Bonfioli\thanks{Lean~4 sources: \texttt{RamaLean/} in the \texttt{rama-notes} repository. Correspondence: \texttt{vicobonfioli@gmail.com}.}}
\date{August 2026}

\begin{document}
\maketitle

\begin{abstract}
This note carries the material removed from \cite{PartII} when that paper was cut to its
spine: the counting refinement and its barriers, the expected non-backtracking polynomial,
the saturation measurements, the coordinate case at rank $b$ with the sharpness of Xus
bound, the account of what a proof would have to contain, and the full formalization record.
Nothing here is needed to read \cite{PartII}.
\end{abstract}


\section{A counting refinement, and its two barriers}

\section{A counting refinement}\label{sec:gapcount}

GAPCOUNT is refuted along with the conjecture it refines, and this section is kept because most
of it does not depend on GAPCOUNT being true. What survives, and is what the section is now for:
the implication GAPCOUNT $\Rightarrow$ localization and the boundary conditions, machine-checked;
BAND, which is a theorem about the abelian shadow and independent of either conjecture; the
settled cases at feedback vertex number one and, conditionally, two; and the two barriers, which
close routes that would otherwise look open. The sweeps are retained as a record of where the
failure is \emph{not} visible. A reader interested only in the biregular case may go straight to
\S\ref{sec:bireg} and the operator form.

Conjecture~\ref{conj:gap} asserts a non-vanishing, which offers no handle. This section
records a strictly stronger statement that equates two integers, the cases it settles, and
the two barriers that close off the obvious routes to the rest. Labels are stated
throughout: what is machine-checked, what is proved modulo a cited classical theorem, and
what is open.

Let $T$ be the universal cover of a finite connected loopless $G$ on $n$ vertices, and
$b=b_1(G)$. Writing $A_G$ for the universal-cover adjacency operator as a matrix over
$C^*_r(F_b)$, a real $x$ lies outside $\operatorname{spec}(T)$ exactly when $xI-A_G$ is
invertible there, and then the trace of its negative spectral projection,
\[
\kappa(x)=\tau_n\bigl(\mathbf 1_{(-\infty,0)}(xI-A_G)\bigr),
\]
is an integer in $[0,n]$, because $K_0(C^*_r(F_b))=\mathbb{Z}\cdot[1]$ by Pimsner--Voiculescu.

\begin{conjecture}[GAPCOUNT]\label{conj:gapcount}
For every $x\notin\operatorname{spec}(T)$, the number of roots of $\mu_G$ strictly above
$x$, with multiplicity, equals $\kappa(x)$.
\end{conjecture}

Conjecture~\ref{conj:gapcount} implies Conjecture~\ref{conj:gap} at $d=1$: $\kappa$ is
constant across a gap while the root count jumps at any root, so a gap contains none.
By the counterexample of \cite{HallCE} the implication now runs the other way and
Conjecture~\ref{conj:gapcount} is false as well, on the same graph: the root $\sqrt5$
crosses inside a gap on which $\kappa$ is constant. We state it as it was, because the
implication and the boundary conditions are what the sweeps below actually test, and
because the counting refinement remains the right object for whichever class of graphs
turns out to satisfy the localization. That
implication, the fact that an integer-valued locally constant function is determined on a
connected gap, and the boundary conditions $\kappa=0$ above $\rho(T)$ and $\kappa=n$ below
$-\rho(T)$, are machine-checked in \texttt{RamaLean/GapLabel.lean}. All the content is in
what happens as a band is crossed.

\subsection{The abelian shadow}

Put a circle variable on each cotree edge and let $A_G(z)$ be the magnetic adjacency matrix,
$B_k=\lambda_k(\mathbb T^b)$ the $k$-th Floquet band, and $\theta_1\ge\dots\ge\theta_n$ the
roots of $\mu_G$.

\begin{proposition}[BAND]\label{prop:band}
$\theta_k\in B_k$ for every $k$. Consequently, for $x\notin\operatorname{spec}(G^{\mathrm{ab}})$
the number of roots above $x$ equals the number of bands above $x$.
\end{proposition}

\begin{proof}
Godsil--Gutman writes $\mu_G$ as the average characteristic polynomial over $\pm1$ edge
signings, and Marcus--Spielman--Srivastava \cite{MSS} make those signings an interlacing
family. The common interlacing lemma brackets the average at every index, not only the
largest, so $\lambda_k(A_s)\le\theta_k\le\lambda_k(A_{s'})$ for two signings. Both are points
of $\mathbb T^b$, and $B_k$ is the continuous image of a connected space, hence an interval.
\end{proof}

Everything downstream of those two classical inputs is machine-checked in
\texttt{RamaLean/BandTheorem.lean}: that bands are order connected
(\texttt{ordConnected\_range\_of\_continuous}), the squeeze
(\texttt{mem\_of\_squeezed}), the count identity (\texttt{countAbove\_eq\_countBands}) and
the recovery of the unicyclic case (\texttt{conj10\_of\_bands\_subset}). The inputs
themselves enter \texttt{band\_of} as hypotheses.

The core of the common interlacing lemma is machine-checked in
\texttt{RamaLean/Interlacing.lean}: a sign argument fixes the parity of a root count and
nothing more, an interlacing confines that count to a window of width two, and parity with
such a window determines the value. What is not formalized, and is cited to \cite{MSS}, is
that the signings have a common interlacing at all.

Proposition~\ref{prop:band} settles Conjecture~\ref{conj:gap} whenever every band lies in
$\operatorname{spec}(T)$, which is exactly $b_1(G)=1$, since $F_1=\mathbb{Z}$ makes the maximal
abelian cover the universal cover. That recovers \cite{PartI} by a different route
and sharpens it to $\theta_k\in B_k$. The hypothesis fails at $b_1\ge2$: for $K_4$ the top
band carries the Perron value $3$ while $\rho(T)=2\sqrt2$.

The counting half of Proposition~\ref{prop:band} does not need \cite{MSS}. If $x$ avoids
every $\lambda_k(z)$ then $\{z:\lambda_k(z)>x\}$ is clopen, so on a connected parameter space
the eigenvalue count above $x$ is the same integer for every $z$; every member of the family
then has the same sign at $x$, so the average does too. Parity follows at once, and parity
upgrades to equality by counting: crossing one component of the spectrum, each component
carries an odd number of roots, and $n$ odd positive integers summing to $n$ are all $1$.
Both steps are machine-checked in \texttt{RamaLean/BandCount.lean}.

\subsection{Feedback vertex number one}

Suppose some vertex $v$ has $F=G-v$ a forest, so every cycle passes through $v$; the first
Betti number is unrestricted, and flowers, theta graphs and every $K_{2,q}$ qualify.

\begin{proposition}\label{prop:fvs1}
Conjecture~\ref{conj:gapcount}, hence Conjecture~\ref{conj:gap} at $d=1$, holds for every
such $G$, modulo Pimsner--Voiculescu and Godsil's path tree.
\end{proposition}

\begin{proof}
Eliminating the preimage of $F$ gives a Schur complement $S_v(x)$ which is a \emph{finite}
sum of group elements, so it lies in the group algebra inside $C^*_r(F_b)$ and not merely in
$L(F_b)$; that distinction is what makes projectionlessness applicable, since $L(F_b)$ is a
$\mathrm{II}_1$ factor full of projections. Haynsworth additivity gives
$\kappa(x)=N_F(x)+\delta(x)$ with $\delta(x)\in\{0,1\}$, and
$\tau(S_v(x))=\mu_G(x)/\mu_{G-v}(x)$, so definiteness fixes the parity of $N_G-N_F$. Godsil's
path tree makes $\mu_{G-v}$ interlace $\mu_G$, confining $N_G-N_F$ to $\{0,1\}$. A window of
width two and a parity determine the value.
\end{proof}

The assembly is \texttt{FeedbackGapCount.gapcount\_fvs\_one\_of}, with the three classical
inputs as explicit hypotheses, and \texttt{conj10\_fvs\_one\_of} composes it with the
skeleton of Section~\ref{sec:gapcount}.

Three remarks on that proof. The combinatorial core, that in an acyclic graph two distinct
neighbours of $v$ cannot be joined without passing through $v$, is machine-checked
(\texttt{FeedbackVertex.no\_bypass}); it is why a walk returning to the same lift must enter
and leave at the same attachment point. The trace formula requires $G$ \emph{simple}: the
argument uses that a lifted attachment point has exactly one neighbouring lift of $v$.
Haynsworth additivity needs no appeal to Sylvester's law in a von Neumann algebra: the block
matrix is congruent to the direct sum of the corner and the Schur complement by explicit
unipotent factors, and scaling the off-diagonal part keeps every point of the path
invertible, so the trace of the negative spectral projection is constant along it. That is
machine-checked in \texttt{RamaLean/Congruence.lean}.

\subsection{Feedback vertex number two, and where it stops}

At $k=2$ the argument breaks in a locatable place: $M_2(C^*_r(F_b))$ is \emph{not}
projectionless, $\operatorname{diag}(1,0)$ being a counterexample, so $S(x)$ can be
indefinite. What survives is the $K$-theory, which keeps $\delta(x)\in\{0,1,2\}$. Two
deletions give a window of width three, so parity leaves a two-way choice between
$\delta=0$ and $\delta=2$, and one further bit is needed. That bit exists.

\begin{proposition}\label{prop:tiebreak}
Let $F=G-v_1-v_2$ and let $a,b_1,b_2,c$ count the roots above $x$ of
$\mu_G,\mu_{G-v_1},\mu_{G-v_2},\mu_F$. If $a-c$ is even then $b_1=b_2$, and
\[
a=c\quad\Longleftrightarrow\quad \mu_F(x)\ \text{and}\ \mu_{G-v_1}(x)+\mu_{G-v_2}(x)\ \text{agree in sign}.
\]
\end{proposition}

\begin{proof}
Vertex-deletion interlacing gives $a-b_i\in\{0,1\}$ and $b_i-c\in\{0,1\}$. If $a=c$ then
$c\le b_i\le a=c$; if $a=c+2$ then $c+1=a-1\le b_i\le c+1$. Either way $b_1=b_2$, so the two
polynomials share a sign at $x$ and so does their sum.
\end{proof}

Proposition~\ref{prop:tiebreak} is machine-checked in \texttt{RamaLean/TieBreak.lean}, and it
needs no hypothesis about $\operatorname{spec}(T)$. It is not enough. The remaining gap is the
parity itself: at $k=1$ it came free from $\tau(S_v)=\mu_G/\mu_F$ and definiteness, while at
$k=2$ the corresponding object is the constant term of the abelianised determinant of a
$2\times2$ matrix over a noncommutative algebra, whose sign has no evident reason to be
$(-1)^{\delta}$. So feedback vertex number two is \emph{reduced, not proved}. The arithmetic
that would close it once the parity is supplied, a window of width three together with a
parity and one tie-break bit, is \texttt{FeedbackTwo.eq\_of\_window3}, and
\texttt{gapcount\_fvs\_two\_sign} discharges the tie-break using
Proposition~\ref{prop:tiebreak}, leaving the parity as the sole remaining hypothesis.

What it reduces to is now quantitative. Since $\mu_G/\mu_F$ is the average over
$\mathbb T^b$ of $\det S(x,z)$, and $\operatorname{sign}\det S(x,z)=(-1)^{\delta_{\mathrm{ab}}(x,z)}$,
writing $I_j$ for the integral of $|\det S|$ over $\{\delta_{\mathrm{ab}}=j\}$ gives
$\mu_G/\mu_F=I_0-I_1+I_2$, and the required sign is exactly the assertion that the parity
class matching $\delta$ dominates the others \emph{in integral}. Measurement
(\texttt{code/inertia\_split.py}) shows the split is lopsided but not degenerate: at two
triangles joined by an edge, in a genuine gap with $\delta=1$, the class
$\delta_{\mathrm{ab}}=1$ occupies $0.987$ of the torus and $\delta_{\mathrm{ab}}=2$ occupies
$0.013$. The integrand therefore does change sign, so no argument asserting that it does not
can work. The decomposition and the criterion are recorded in
\texttt{RamaLean/InertiaSplit.lean}, with a witness that measure alone does not decide it.

\subsection{Two barriers}

Both close off routes that look available.

\begin{proposition}\label{prop:nohom}
For $b\ge2$ there is no unital $*$-homomorphism $C^*_r(F_b)\to C(\mathbb T^b)$ carrying the
generators to the coordinates.
\end{proposition}

\begin{proof}
Such a map would carry units to units, so $S(x)$ invertible would force $S(x,z)$ invertible
for every $z$, hence $\det(xI-A_G(z))=\mu_F(x)\det S(x,z)$ nonzero for every $z$, hence
$\operatorname{spec}(G^{\mathrm{ab}})\subseteq\operatorname{spec}(T)\cup\{\mu_F=0\}$. For
$K_4$ with a two-element feedback set that fails at $x=3$: the Perron value lies in
$\operatorname{spec}(G^{\mathrm{ab}})$, while $\rho(T)=2\sqrt2<3$ and $\mu_{K_2}(3)=8$.
\end{proof}

The refutation shape is machine-checked in \texttt{RamaLean/NoQuotient.lean}. The underlying
reason is that a quotient descends to reduced $C^*$-algebras when the \emph{kernel} is
amenable, and $[F_b,F_b]$ is free of infinite rank.

The second barrier is narrower than it first appeared. Every argument that runs on the torus family needs $x$ outside
$\operatorname{spec}(G^{\mathrm{ab}})$, and on the smallest examples that fails: for $K_4$,
two triangles joined by an edge, the theta graph and $K_4$ with pendants, the internal gaps
of $\operatorname{spec}(T)$ are swallowed by bands of the abelian cover. It would be wrong to
read that as a general obstruction. A certified search over two growing families
(\texttt{code/torus\_gg.py}) finds the opposite behaviour as soon as the graphs are larger:
ten of fourteen gap midpoints lie outside $\operatorname{spec}(G^{\mathrm{ab}})$, with a
margin that grows with the number of vertices: for two $c$-cycles joined by an edge it runs
$1.18,\,7.52,\,21.50,\,51.35$ as $c$ goes from four to seven. So the torus family is vacuous
on the smallest graphs and informative as they grow, the reverse of what the first
computation suggested, and the localization of \cite{PartI} settles
Conjecture~\ref{conj:gap} outright at every such point, for every first Betti number and with
no feedback vertex hypothesis.

Nor is that settlement subsumed by Heilmann--Lieb. Those graphs have maximum degree three, so
Heilmann--Lieb excludes roots only outside $[-2\sqrt2,2\sqrt2]$, while the certified points
sit at $\pm1.93,\pm1.99,\pm2.04,\pm2.06$, all interior. The certification cannot be a grid
minimum, which proves nothing when the zero set of $\det(xI-A_G(z))$ is a curve: that
determinant has degree one in each $z_e$, so its $3^{b_1}$ Fourier coefficients are recovered
exactly by a three-point transform, they bound the gradient, and a grid search together with
that Lipschitz constant certifies a positive lower bound, while a sign change on the grid
certifies the opposite verdict.

\subsection{Evidence}

\texttt{code/universal\_cover.py} computes $\operatorname{spec}(T)$ and $\kappa$ for an
arbitrary finite graph, by the non-backtracking cavity equations on the directed edges. Its
regression check is $K_4$ with a pendant at each vertex, where
$\operatorname{spec}(T)=\pm[\sqrt3-\sqrt2,\sqrt3+\sqrt2]$ in closed form; the solver returns
$[-3.1456,-0.3163]$ and $[0.3184,3.1477]$. Runs whose density fails to integrate to one are
discarded, since point spectrum makes the cavity fixed point singular and a diverged
$\kappa$ reads exactly like a violated Conjecture~\ref{conj:gapcount}. That the cavity map
preserves the lower half plane, so the iteration is total and the density non-negative, is
machine-checked in \texttt{RamaLean/Cavity.lean}; convergence is not claimed.

A Rust sweep (\texttt{code/covercheck}) applies this to every connected graph of feedback
vertex number at most two on up to seven vertices, computing the root count by
Budan--Fourier, which is exact for a real-rooted polynomial, and $\kappa$ by the cavity
equations. Over $55\,131$ graphs and $99\,487$ gap points there is no failure of
Conjecture~\ref{conj:gap} and none of Conjecture~\ref{conj:gapcount}, no run is discarded by
the validity gate, and $\kappa$ is never further than $0.205$ from an integer.
That range could not have found the counterexample and should not be read as having looked:
Hall's graph has $41$ vertices and feedback vertex number five, so it lies outside the sweep in
both parameters. What the sweep establishes is that the failure is not visible at small size,
which is why it went unnoticed, and not that it is rare.

The measurement also shows what an estimate would have to exploit. The wrong-parity class
is not merely small in measure, it is suppressed quadratically in integral: at two triangles
joined by an edge, in the gap at $x=-1.767$, it occupies $0.013$ of the torus and
contributes $0.0007$ of the integral of $|\det S|$, a margin of $0.9987$. The reason is
structural: the wrong-parity region is a shell around $\{z:\det S(x,z)=0\}$, where one
eigenvalue has just crossed and $|\det S|$ is correspondingly small. Small measure and small
integrand together, which is why the competition is not close.

Scanning a whole gap rather than one point confirms the exponent. Across both internal gaps
of two triangles joined by an edge, $I_{\mathrm{wrong}}/m_{\mathrm{wrong}}^2$ stays between
$2.4$ and $5.2$ and the margin never falls below $0.977$
(\texttt{code/margin\_scan.py}); there is no blow-up at the gap edges, where the
wrong-parity class is largest and the ratio is smallest. The chain that would turn this into
an estimate is machine-checked in \texttt{RamaLean/ShellBound.lean}. Its first step is an
identity rather than a bound: for a $2\times2$ Hermitian matrix the operator norm is the
larger eigenvalue in absolute value, so
$|\det S|=\lVert S\rVert\cdot\min(|\lambda_1|,|\lambda_2|)$
(\texttt{abs\_mul\_eq\_max\_mul\_min}), and on the shell only the crossing eigenvalue
matters. Weyl makes that eigenvalue Lipschitz in the phase, giving
$|\det S(z)|\le ML\operatorname{dist}(z,\partial)$ and hence
$I_{\mathrm{wrong}}\le m_{\mathrm{wrong}}\,ML\operatorname{reach}$, where
$\operatorname{reach}=\sup_{z\in A}\operatorname{dist}(z,\partial A)$ is the inradius
(\texttt{det\_le\_on\_shell}). The relevant quantity is the inradius and not the diameter: a
shell wrapping the torus has diameter of order one however thin it is.

The inradius is then bounded for free. A region of inradius $r$ contains a ball of radius
$r$, so $c_br^b\le m$ and $\operatorname{reach}\le(m/c_b)^{1/b}$, giving
\[
   I_{\mathrm{wrong}}\ \le\ ML\,c_b^{-1/b}\,m_{\mathrm{wrong}}^{\,1+1/b},
\]
which at $b=2$ is $m^{3/2}$ (\texttt{reach\_le\_sqrt},
\texttt{integral\_le\_three\_halves}). The upper half of the estimate is therefore
unconditional. The region is not a thin shell but a blob, with $m$ of order
$\operatorname{reach}^2$: measurement gives $\operatorname{reach}/m$ running $2.60$, $3.10$,
$4.63$, $5.94$, $9.80$ as $m$ falls from $0.052$ to $0.0053$. The exponent $3/2$ is accordingly
weaker than the $m^2$ the numerics suggest, and stronger than anything needed. The lower half, a bound $I_{\mathrm{right}}\ge c>0$, is open, and
its natural route is precisely what Proposition~\ref{prop:nohom} rules out.

The obvious substitute is circular, and it is worth saying why. Splitting the average by
inertia class gives $\mu_G/\mu_F=I_0-I_1+I_2$ \emph{identically}, so
$I_{\mathrm{right}}-I_{\mathrm{wrong}}=|\mu_G/\mu_F|$ and bounding $I_{\mathrm{right}}$
below by $|\mu_G/\mu_F|$ assumes exactly the sign at issue. No independent lower bound on
$|\mu_G(x)|$ is available inside a gap, since Heilmann--Lieb gives root-free intervals only
outside $[-\rho,\rho]$, where the conjecture is already known.

A substitute that is not circular is the geometric mean. Jensen gives
\[
   I_{\mathrm{right}}+I_{\mathrm{wrong}}\;=\;\int_{\mathbb T^b}|\det S|\;\ge\;
   \exp\Bigl(\int_{\mathbb T^b}\log|\det S|\Bigr)\;=:\;\Delta ,
\]
a Mahler measure of the abelian cover, which owes nothing to the sign of $\mu_G$. Hence
$I_{\mathrm{right}}\ge\Delta-I_{\mathrm{wrong}}$, and domination holds as soon as
$\Delta>2I_{\mathrm{wrong}}$. Combined with the unconditional upper half this reduces
feedback vertex number two to a single inequality, $\Delta>2ML c^{-1/2}m^{3/2}$
(\texttt{RamaLean/JensenRoute.lean}). The geometric mean, which appeared above as the source
of the obstruction, is here the tool.

Measured across both internal gaps of two triangles joined by an edge, $\Delta$ lies between
$0.759$ and $0.823$ while $I_{\mathrm{wrong}}$ runs from $10^{-4}$ to $10^{-2}$, so
$\Delta/2I_{\mathrm{wrong}}$ runs between $41$ and $3927$
(\texttt{code/jensen\_route.py}); the same script verifies the circularity identity to
$10^{-15}$. A sweep over $1475$ gap points of graphs of feedback vertex number at most two
with $2\le b\le4$ records no failure of $\Delta>2I_{\mathrm{wrong}}$ and a worst ratio of
$298$, and the ratio does not degrade with $b$: its minimum is $298$ at $b=2$, $2070$ at
$b=3$ and $1740$ at $b=4$ (\texttt{code/jensen\_sweep.py}). What is a Mahler measure over
$|\mu_F(x)|$ is recorded in \texttt{RamaLean/MahlerRoute.lean}.

\subsection{Why the feedback vertex set should be removed}\label{sec:parity}

The inequality $\Delta>2MLc^{-1/2}m^{3/2}$ is not the right target, and the reason is an
artefact of the Schur complement rather than a feature of the conjecture. Both $M=\sup\|S\|$
and the Lipschitz constant $L$ scale like $1/\mu_F(x)$, while $\Delta$ scales like
$1/\mu_F(x)$; so as $x$ approaches a root of $\mu_F$ inside a gap the left side grows one
whole factor of $\mu_F$ faster than the right, and the inequality fails at points where
nothing is wrong. Roots of $\mu_F$ are not excluded from gaps of $\operatorname{spec}(T)$.

The Schur complement is a computational device and can be dropped. With
$P_x(z)=\det(xI-A_G(z))=\prod_k(x-\lambda_k(z))$ and $N(z)=\#\{k:\lambda_k(z)>x\}$, the sign
of $P_x$ is $(-1)^{N(z)}$, and the localization of \cite{PartI} gives
\[
   \mu_G(x)=\int_{\mathbb T^b}P_x
   =I_{\mathrm{even}}-I_{\mathrm{odd}},\qquad
   I_p=\int_{\{N\equiv p\ (2)\}}|P_x| ,
\]
with no feedback vertex hypothesis, no restriction on $b$, and no $\mu_F$. It also makes the
Lipschitz constant explicit instead of measured, and better than the obvious bound. Crudely,
$\partial A/\partial\theta_j$ has a single entry of modulus one together with its conjugate,
so its operator norm is $1$ and $\|A(z)-A(w)\|\le\sum_j|\Delta\theta_j|$, which costs a factor
$b$. First-order perturbation theory does better: $\partial\lambda_k/\partial\theta_j$ is the
quadratic form of that one edge against a unit eigenvector $v$, so it is at most
$|v_{u_j}|^2+|v_{v_j}|^2$, and summing over cotree edges charges each vertex once per incident
cotree edge. Hence
\[
   \sum_j\Bigl|\frac{\partial\lambda_k}{\partial\theta_j}\Bigr|\;\le\;D ,
\]
with $D$ the largest number of cotree edges at a vertex. $D\le2$ whenever the cotree edges can
be chosen vertex-disjoint and $D\le\Delta(G)$ always, so $D$ does not grow with the graph while
$b$ does. The per-edge estimate (\texttt{BandLipschitz.quad\_form\_bound}), the double count
(\texttt{sum\_deg\_bound}) and the assembly (\texttt{grad\_l1\_bound}) are machine-checked;
first-order perturbation theory is not in Mathlib and enters \texttt{grad\_l1\_bound} as a
hypothesis in the shape the argument consumes.

Let $\kappa$ be the number of bands whose range contains $x$. If $\kappa=0$ the localization
settles $x$. Over $149\,764$ residue points from $40\,146$ graphs on up to eight vertices,
$|K|=1$ at $140\,661$ and $|K|=2$ at $9103$, and no point has $|K|\ge3$
(\texttt{code/jsweep}). How often $|K|=0$ happens depends sharply on $b$: $80.2\%$ of $486$
midpoints
with $b=2$ lie outside $\operatorname{spec}(G^{\mathrm{ab}})$, against $3.4\%$ of $493$
midpoints with $b=3$, where a further $69.8\%$ are certified inside and $26.8\%$ undecided,
so at most $30\%$ of the $b=3$ points can be settled this way
(\texttt{code/residue\_sweep.py}, \texttt{code/residue\_b3.py}). More torus directions fill
the gaps in. Over $5971$ residue points, with $\kappa=1$ at $5397$ and $\kappa=2$ at $574$
and no point with $\kappa\ge3$, the minority-parity region has measure at most
$0.0674$ and
$\min(I_{\mathrm{even}},I_{\mathrm{odd}})/\max(I_{\mathrm{even}},I_{\mathrm{odd}})$ is at
most $0.0169$ (\texttt{code/minority.py}), so the two classes are not close to cancelling.

\begin{proposition}\label{prop:wmean}
Suppose exactly one band function $\lambda_{k_0}$ attains $x$ on $\mathbb T^b$. Then
$Q=\prod_{k\ne k_0}(x-\lambda_k)$ has constant sign $s$, and
\[
   \mu_G(x)=s\int_{\mathbb T^b}\bigl(x-\lambda_{k_0}(z)\bigr)\,|Q(z)|\,dz .
\]
Consequently $\mu_G(x)=0$ if and only if $x$ is the $|Q|$-weighted mean of $\lambda_{k_0}$.
\end{proposition}

\begin{proof}
The factors other than $k_0$ never vanish and $\mathbb T^b$ is connected, so $Q$ has one sign.
The integral is affine in $x$, so it vanishes exactly when
$x\int|Q|=\int\lambda_{k_0}|Q|$.
\end{proof}

Machine-checked as \texttt{ParitySplit.zero\_iff\_weighted\_mean} and
\texttt{ne\_zero\_iff\_ne\_weighted\_mean}, with the identity verified to $10^{-15}$ in
\texttt{code/kap.py}. A weighted mean lies between the bounds of the function averaged
(\texttt{mean\_mem\_Icc}), so an $x$ outside the range of the crossing band cannot be that
mean: \texttt{ne\_zero\_of\_outside\_band} recovers the localization in this case with no
estimate at all. Conjecture~\ref{conj:gap} at a point with $\kappa=1$ becomes the statement
that the $|Q|$-weighted mean of the crossing band lies in $\operatorname{spec}(T)$. Measured
means are $-2.149,-1.483,1.483,2.143,2.150$ on gaps whose interiors sit near $\pm1.77$, so
they are far from $x$ and inside $\operatorname{spec}(T)$.

Two caveats keep this in proportion. First, over $4900$ gap midpoints $\kappa=1$ at $2206$
residue points and $\kappa\ge2$ at $183$, so Proposition~\ref{prop:wmean} covers $92.3\%$ of
the residue and not all of it. Second, the statement it reduces to,
\[
   \text{the } |Q|\text{-weighted mean of the crossing band lies in }\operatorname{spec}(T),
\]
is a conjecture and is labelled as one; what can be said is that a deliberate attempt to
break it failed. Over all $2206$ points with $\kappa=1$ the weighted mean lies in
$\operatorname{spec}(T)$ every time and is never closer than $0.167$ to $x$
(\texttt{code/wmean\_test.py}), and a single coincidence there would have refuted
Conjecture~\ref{conj:gap} outright.

\subsection{Separating the band geometry from the weight}\label{sec:split}

Proposition~\ref{prop:wmean} mixes two unrelated things: where the crossing band sits
relative to $x$, and how much the remaining factors vary. Separating them removes the
restriction to one crossing band and gives a criterion whose two halves can be attacked
independently.

Let $K=\{k:x\in B_k\}$ be the crossing set and put
\[
   \pi(z)=\prod_{k\in K}\bigl(x-\lambda_k(z)\bigr),\qquad
   w(z)=\Bigl|\prod_{k\notin K}\bigl(x-\lambda_k(z)\bigr)\Bigr| .
\]
The factors outside $K$ never vanish, so their product has a constant sign $s$ and
$\mu_G(x)=s\int\pi w$.

\begin{proposition}\label{prop:split}
Write $\pi=\pi_+-\pi_-$ and $J_\pm=\int_{\mathbb T^b}\pi_\pm$, and suppose
$0<q_{\mathrm{lo}}\le w\le q_{\mathrm{hi}}$. If
\[
   \frac{J_+}{J_-}\;>\;\Lambda:=\frac{q_{\mathrm{hi}}}{q_{\mathrm{lo}}} ,
\]
then $\mu_G(x)\neq0$. No hypothesis on $|K|$ is used.
\end{proposition}

\begin{proof}
$\int\pi w=\int\pi_+w-\int\pi_-w\ge q_{\mathrm{lo}}J_+-q_{\mathrm{hi}}J_->0$.
\end{proof}

Machine-checked as \texttt{CrossingSplit.ne\_zero\_of\_split}, with the ratio form as
\texttt{crit\_of\_ratio}. The point is that $J_\pm$ involve only the crossing bands and
$\Lambda$ only the rest, so a bound on either half is useful on its own; the weighted mean of
Proposition~\ref{prop:wmean} admits no such division. It is also not the wasteful estimate:
bounding the whole integrand by $\sup w$ loses the four orders of magnitude that
$\sup w/\inf w$ can reach, whereas here $\inf w$ is applied to the majority term and $\sup w$
only to the minority term.

The criterion is not always satisfied. Over $149\,764$ residue points drawn from $40\,146$
graphs on up to eight vertices it holds $149\,511$ times and fails $253$ times, with worst
margin $0.0616$ (\texttt{code/jsweep}). A sample of $208$ points had shown it holding every
time, which was too small to see the failures.

Two facts sharpen what is left. Every failure has exactly one crossing band: at $|K|=2$ the
criterion holds at all $9103$ points, and no point with $|K|\ge3$ occurs in the corpus at all.
And Conjecture~\ref{conj:gap} is untouched at the failing points, the true ratio
$\min(I_+,I_-)/\max(I_+,I_-)$ never exceeding $0.0233$ over the whole corpus, so the
conclusion holds comfortably where this sufficient condition does not see it. The worst case
is $n=8$, $b=3$, $|K|=1$, $x=2.0607$ on the graph with edges
$01,02,03,04,14,25,36,37,56,57$.

What the criterion buys is therefore a reduction covering $99.83\%$ of the residue, together
with a target that has structure: $J_+/J_-$ is large because $x$ sits near an edge of the
crossing band, the minority measure never exceeding $0.0674$, and $\Lambda$ is a product of
band-distance ratios over the non-crossing bands. Bounding either one is a self-contained
problem, and the $253$ exceptions are a small explicit list to be understood rather than a
diffuse obstruction.

The weight can be removed from the criterion entirely. A product is bounded above by the
product of the upper bounds of its factors and below by the product of the lower bounds, and
for a band that does not cross $x$ the range $B_k=[\ell_k,h_k]$ lies on one side, so the
corresponding factor is exactly $1+\operatorname{width}(B_k)/\operatorname{dist}(x,B_k)$.
Hence
\[
   \Lambda\;\le\;\prod_{k\notin K}\Bigl(1+\frac{\operatorname{width}(B_k)}
   {\operatorname{dist}(x,B_k)}\Bigr),
\]
machine-checked as \texttt{WeightBound.crit\_of\_band\_geometry} with the per-band identity as
\texttt{factor\_below} and \texttt{factor\_above}. The criterion then reads entirely in the
band ranges of the abelian cover, in the same currency as $J_+/J_-$.

That bound is not sufficient on its own, and the failure is measured. The only thing discarded
is the correlation between bands, since $\sup\prod\le\prod\sup$ is an equality only when the
extrema align, and that costs a median factor $4.73$ and up to $28.7$. Over $208$ residue
points the criterion with the true $\Lambda$ holds every time, worst margin $2.149$, while the
criterion with the product in place of $\Lambda$ holds $207$ times and fails once, at
$J_+/J_-=1213$ against a product of $2298$. The failing point is the one where $x$ sits
closest to the edge of the crossing band, minority measure $0.0171$, which is not a
coincidence: proximity to a band edge makes $J_+/J_-$ large and also tends to bring a second
band close to $x$, which makes the product large. Decoupling the two halves therefore costs
most exactly where the margin is thinnest.

Difficulty is conserved and it is worth saying where it went. A weighted mean lies in the
convex hull of the range of $\lambda_{k_0}$, which is the band containing $x$, and that band
is not inside $\operatorname{spec}(T)$ precisely because $x$ is in it and outside
$\operatorname{spec}(T)$. So the containment is not automatic. What has changed is the shape
of the remaining work: one scalar in one interval in place of two integrals of comparable
size. Two crude routes to it are already closed. Replacing $|Q|$ by its extremes fails
because $\sup|Q|/\inf|Q|$ reaches $5.4\cdot10^{4}$ on the same examples, which is far more
than the small minority measure buys. And deforming the polytorus radii, which leaves the
constant Fourier coefficient unchanged and so preserves $\mu_G(x)$, cannot help either: at
radius one the image of $P_x$ lies on a line and spans an arc of exactly $\pi$, while any
other radius makes $A_G$ non-Hermitian and opens the image into two dimensions, so the arc
only grows and the half-plane criterion is never reached (\texttt{code/halfplane.py}).


\section{The expected non-backtracking polynomial, and how close the bound comes}

\section{An expected non-backtracking polynomial}

\begin{lemma}[Chernyak--Chertkov; Watanabe--Fukumizu]\label{lem:gauge}
Let $G$ be any graph, $A_s$ its adjacency matrix with edges signed by
$s\in\{\pm1\}^{E}$, and $D$ the degree matrix. Then, with $s$ uniform,
\[
\mathbb E_s\,\det\bigl(I+u^2(D-I)-uA_s\bigr)
=\sum_{M}(-u^2)^{|M|}\prod_{v\notin V(M)}\bigl(1+(d_v-1)u^2\bigr),
\]
the sum being over all matchings $M$ of $G$. If $G$ is $(a,b)$-biregular
bipartite with $P=a-1$ and $Q=b-1$, the right side is $\mu_G(\lambda)$ up to
normalization at $\lambda=\sqrt{(1+Pu^2)(1+Qu^2)}/u$, and then, with $t=u^2$,
\begin{equation}\label{eq:tsub}
\lambda^2-c=\frac1t+mt .
\end{equation}
\end{lemma}

The identity is not new and we claim none of it. Chernyak and Chertkov
\cite[Eqs.~(2),~(6)]{CC} prove that for an arbitrary matrix $H$ the average of
$\det H(\sigma)$ over $\mathbb Z_2$ edge gaugings equals a monomer--dimer sum with
weights $w_a=H_{aa}$ and $w_{ab}=-H_{ab}H_{ba}$; taking $H=I+u^2(D-I)-uA$ gives
the display, since then $w_a=1+(d_a-1)u^2$ and $w_{ab}=-u^2$. Equivalently, the
averaging step alone is the Godsil--Gutman expansion \cite{GG}, valid for any
diagonal matrix in place of $I+u^2(D-I)$, and the right-hand side with those
degree-dependent weights is the monomer--dimer partition function
$\Xi_G(-u^2,\lambda)$, $\lambda_v=1+(d_v-1)u^2$, of Watanabe and Fukumizu
\cite[Thm.~5.5]{WF}; on a $(q+1)$-regular graph they reduce it to
$u^{|V|}\mu_G(1/u+qu)$ \cite[Eq.~(5.3)]{WF}, attributing the regular case to
Nagle, an attribution we take from \cite{WF}, not having consulted the
original, and record the consequent band statement, that the roots then lie on
the
circle of radius $1/q$. Only \eqref{eq:tsub}, the biregular substitution, is not
in those sources.

What we do wish to point out is how the left-hand side should be read. By the
Ihara--Bass formula for signed graphs, $\mathbb E_s\det(I-uW_s)$ is
$(1-u^2)^{|E|-|V|}$ times the left side above, so the polynomial in
Lemma~\ref{lem:gauge} is an \emph{expected non-backtracking characteristic
polynomial}. This reading is absent from the sources just cited, and not by
oversight in an obvious direction: \cite{CC} obtain the average but explicitly set
the Ihara spectral parameter to zero, and \cite{WF} display the Bass matrix
$I-uA_G+u^2(D_G-I)$ next to their Theorem 5.9, a sum over deletions of
vertex-disjoint cycles, and remark in print that its relation to the graph zeta
function is unknown. The sign average is exactly what annihilates those cycle
terms. The closest antecedent we have found is Watanabe and Chertkov
\cite[Thm.~4]{WC}, who do average a non-backtracking determinant over $\pm1$ edge
signs and convert it by Ihara--Bass, but for the permanent, that is for perfect
matchings on $K_{N,N}$, with no degree-dependent vertex weights.

The band condition has a clean reading on that polynomial, and it sharpens a
published bound. By \eqref{eq:tsub}, a root $\beta$ of $\omega_G$ coming from a
nonzero root $\lambda$ of $\mu_G$ satisfies $m\beta^2-\rho\beta+1=0$ with
$\rho=\lambda^2-c$, so the two roots of that quadratic have product $1/m$; they are
complex conjugates, hence both of modulus $m^{-1/2}$, exactly when
$|\rho|\le2\sqrt m$, and otherwise real with one inside and one outside that circle.
Watanabe and Fukumizu prove \cite[Cor.~5.8]{WF} that every root of $\omega_G$ lies
in the annulus $1/(d_M-1)\le|\beta|\le1/(d_m-1)$, which for an $(a,b)$-biregular
graph reads $1/(a-1)\le|\beta|\le1/(b-1)$. Conjecture~\ref{conj:gap} at $d=1$
asserts that all roots except those coming from $\lambda=0$ lie on the single
circle $|\beta|=m^{-1/2}$, the geometric mean of the two radii they bound by. The
roots at $\lambda=0$ are $\beta=-1/(a-1)$ and $\beta=-1/(b-1)$, so they sit at the
two ends of that annulus, which is why it is attained and cannot be narrowed
without excluding them.

Since $1/t+mt=2\sqrt m\cos\varphi$ exactly when $|t|=1/\sqrt m$:

\begin{corollary}\label{cor:ramanujan}
For $(a,b)$-biregular bipartite $G$, Conjecture~\ref{conj:gap} at $d=1$ holds for
$G$ if and only if every root $u$ of the expected non-backtracking characteristic
polynomial of $G$, other than the trivial ones, satisfies $|u|=m^{-1/4}$.
\end{corollary}

So the conjecture is a Ramanujan condition after all, but for an \emph{expected}
polynomial rather than for a spectrum. That is forced. If $\mu_G$ were, up to the
substitution and prefactor above, the determinant of a single locally described
operator on $G$, its coefficients could be recovered by interpolation from
polynomially many determinant evaluations; but those coefficients are the
matching numbers $m_k$, whose computation is $\#P$-hard even for planar graphs of
bounded degree \cite{Je}. No unaveraged Bass-type identity for the matching
polynomial can therefore exist unless $\mathrm{FP}=\#\mathrm P$. This is also why
the roots of $\mu_G$ are not the spectrum of any operator on $G$, and hence why
the available technology for such statements is the method of interlacing
families \cite{MSS}, which works with expected polynomials, rather than a
spectral argument.

\section{How close the bound comes to being attained}

Proposition~\ref{prop:moments} also explains why the interval is not filled. The
mean square of the roots is $\bigl[(b-2)^2+a(b-1)\bigr]$, so the fraction of the
band $[-2\sqrt m,2\sqrt m]$ occupied in mean square is
\[
\frac{(b-2)^2+a(b-1)}{4(a-1)(b-1)},
\]
which is $7/16$ for $(a,b)=(3,3)$, $1/2$ for $(2,2)$, $3/8$ for $(4,3)$, and
tends to $1/4$ as $a\to\infty$. The bulk of the spectrum sits well inside the
band; only the extremes are at issue, which is why moments cannot decide the
question.

The extremes behave as Theorem~\ref{thm:subcoef} predicts. Writing the least and
greatest roots as percentages of the two band edges:
\begin{center}
\begin{tabular}{lccc}
\hline
$G$ & girth & least root & greatest root\\
\hline
$K_{3,3}$ & 4 & $89.6\%$ & $57.2\%$\\
$K_{4,4}$ & 4 & $94.6\%$ & $56.6\%$\\
circulant $6+6$ & 4 & $96.4\%$ & $74.2\%$\\
Heawood & 6 & $97.5\%$ & $79.9\%$\\
Pappus & 6 & $98.6\%$ & $83.4\%$\\
\hline
$C_{12}$ & 6 & $96.6\%$ & $96.6\%$\\
$3C_4$ & 4 & $70.7\%$ & $70.7\%$\\
\hline
\end{tabular}
\end{center}
The first block has $b=3$, and the least root is consistently far closer to its
edge than the greatest is to its, the roots are displaced downwards by $b-2$,
exactly as Theorem~\ref{thm:subcoef} requires. The second block has $b=2$, the
displacement vanishes, and the two ends are used equally. That is what the table
shows, and it is all it shows.

It does not bear on sharpness, for a reason worth stating because it is easy to
miss. Every graph in the table has $a=b$, and there $c=2(a-1)=2\sqrt m$, so
$\rho=y-c\ge-c=-2\sqrt m$ holds automatically from $y=x^2\ge0$: the lower bound is
an identity, not a theorem. This is the triviality that Song, Fan and Miao note
for $d=r$ \cite{SFM}. The percentages above therefore measure only how close the
least root of $\mu_G$ comes to $0$, and $100\%$ would require $\mu_G(0)=0$, that
is a graph with no perfect matching, which by K\H onig's theorem no regular
bipartite graph is. The conjecture has content exactly when $a\ne b$, where
$c>2\sqrt m$ strictly by the arithmetic--geometric mean inequality.

Where it has content the bound holds with room to spare, and is not approached.
For $(4,2)$-biregular graphs the conjecture asserts $y\ge c-2\sqrt m=4-2\sqrt3
=0.5359$; a girth-$4$ family plateaus at $y_{\min}=1.1716$ and a girth-$8$ family
at $0.7392$. In the trivial regime the same plateauing is visible, and there
the limit is an algebraic number one can name. For a cubic bipartite circulant on
$n+n$ vertices with connection set $S$, the matching counts satisfy
$f_n(y)=\operatorname{tr}A(y)^n$ for a transfer matrix $A(y)$ with entries linear
in $y$, so by the theorem of Beraha, Kahane and Weiss the roots accumulate where
the two dominant eigenvalues of $A(y)$ have equal modulus. That happens by a
genuine collision of two real eigenvalues, below $y^*$ the pair is real and
distinct, above it a complex-conjugate pair, so the locus is cut out by
$\operatorname{disc}_\lambda\chi_{A(y)}=0$ and $y^*$ is the left edge of the
accumulation band. For $S=\{0,1,2\}$, girth $4$,
\[
\operatorname{disc}_\lambda\chi=4y^2\bigl(y^4-12y^3+46y^2-84y+5\bigr),
\]
and $y^*$ is the least positive root of that quartic, which is irreducible over
$\mathbb Q$; then $y^*=0.0615663466041589\ldots$ and the limit is
$100-25y^*=98.4608413349\ldots\%$ of the lower edge. For $S=\{0,1,3\}$, girth
$6$, the discriminant is $y^3F(y)$ with $F$ irreducible of degree $17$, giving
$y^*=0.0059255959422562\ldots$ and $99.8518601014\ldots\%$. Both stop strictly
short of $100\%$. The finite-$n$ values approach these from above with a $1/n^2$
band-edge law, and $n^2\bigl(y_{\min}(n)-y^*\bigr)$ is constant to five digits at
$n=200,300,500$, which identifies the extrapolated plateau with the algebraic
number.

One half of the limit is certified exactly, without appeal to
Beraha--Kahane--Weiss. On the whole of $[0,y^*)$ the largest eigenvalue modulus of
$A(y)$ is attained by exactly one eigenvalue, which is real, simple and negative:
the discriminant has no zero there, so real branches stay real; no two eigenvalues
are negatives of one another, by a resultant of the even and odd parts of $\chi$
which has no zero on $(0,y^*]$; and for the octic a separating radius $5/4$, which
no eigenvalue ever attains, together with one exact Schur count, pins the number of
eigenvalues outside it at two. Hence $|f_n|\ge|\lambda_1|^n(1-(N-1)(1+\delta)^{-n})$
on any compact subinterval, giving $\liminf_ny_{\min}(n)\ge y^*$. The reverse
inequality does use Beraha--Kahane--Weiss: just above $y^*$ the two branches form a
conjugate pair of equal and maximal modulus, and one needs that theorem to conclude
that the zeros accumulate there. This is what one should expect:
a family with a fixed local structure has a fixed Benjamini--Schramm limit, which
is not a tree, so its matching measure retains a gap at $0$. Non-attainment of
this kind was already visible to Heilmann and Lieb, who tabulate the largest zero
against their own bound for six lattices, $11.24$ against $12$, $17.86$ against
$20$, $41.32$ against $44$, and remark that equality ``cannot be expected to be
true in general'' \cite[Table 1 and Remark 4.3]{HL}. Girth does help
monotonically, at $30$ vertices, girth $4$, $6$ and $8$ give $98.166\%$,
$99.311\%$ and $99.410\%$, but no family of bounded girth can approach the
edge, and every cubic bipartite circulant has girth at most $6$, since the closed
hexagon $0,s_1,s_1-s_2,s_1-s_2+s_3,s_3-s_2,s_3$ is always present. We therefore
make no claim that $2\sqrt m$ is attained or approached. The roots of these
circulants do satisfy \eqref{eq:twosided}, with maxima $4.06,4.48,6.26$ against
bounds $4.90,5.66,6.93$.


\section{The coordinate case at rank $b$, and the sharpness of Xus bound}

\subsection{The coordinate case at rank \texorpdfstring{$b$}{b}, and the sharpness of Xu's
constant}\label{sec:xusharp}

Theorem~\ref{thm:matchingform} identifies the recentred mixed characteristic polynomial with the
matching polynomial in the commuting case at $b=2$. That identification survives to every rank,
and the object it produces is again an ordinary matching polynomial, on a different graph. The
consequence is that Xu's constant is not merely an upper bound at $b\ge3$: it is exactly the
truth, and the class attains it.

Let $H$ be an $a$-regular $b$-uniform hypergraph on $n$ vertices with $q$ hyperedges, and let
$P_e$ be the coordinate projection onto $\operatorname{span}\{e_v:v\in e\}$. These have rank $b$
and satisfy $\sum_eP_e=aI$ exactly, so they are an admissible tight family. Write $I(H)$ for the
incidence bipartite graph, with one side the $n$ vertices and the other the $q$ hyperedges; it is
$(a,b)$-biregular.

\begin{proposition}\label{prop:bridge}
With $\mu$ the mixed characteristic polynomial of $\{P_e\}$,
\[
\mu(x)=\sum_{k\ge0}(-1)^k\,m_k\bigl(I(H)\bigr)\,x^{\,n-k},
\qquad\text{equivalently}\qquad
\mu(t^2)=t^{\,n-q}\,\mu_{I(H)}(t),
\]
where $m_k$ is the number of $k$-matchings. In particular the positive roots of $\mu$ are exactly
the squares of the positive roots of $\mu_{I(H)}$.
\end{proposition}

\begin{proof}
The matrix $xI+\sum_ez_eP_e$ is diagonal, with $v$-th entry $x+\sum_{e\ni v}z_e$, so
$\det=\prod_v\bigl(x+\sum_{e\ni v}z_e\bigr)$. Each factor is affine in every $z_e$, so for
$S\subseteq[q]$ the derivative $\partial_S\det$ evaluated at $z=0$ is obtained by choosing, for
each $e\in S$, one factor to differentiate, that is one vertex $v\in e$; two hyperedges choosing
the same vertex differentiate one affine factor twice and contribute $0$. Hence $\partial_S\det$
at $z=0$ counts the injective maps $S\to V$ with $e\mapsto v\in e$, each contributing
$x^{\,n-|S|}$. Applying $\prod_e(1-\partial_{z_e})$ and grouping by $|S|=k$, the coefficient is
the number of systems of distinct representatives of $k$ hyperedges, which is the number of
$k$-matchings of $I(H)$ saturating $k$ hyperedge-nodes. The second display follows from
$\mu_{I(H)}(t)=\sum_k(-1)^km_kt^{\,n+q-2k}$ on substituting $x=t^2$.
\end{proof}

\begin{corollary}\label{cor:xucoord}
For every coordinate family as above, at every $b\ge2$,
\[
\operatorname{maxroot}\mu\ \le\ \bigl(\sqrt{a-1}+\sqrt{b-1}\bigr)^2 .
\]
\end{corollary}

\begin{proof}
$I(H)$ is $(a,b)$-biregular, so its universal cover is the $(a,b)$-biregular tree, of spectral
radius $\rho=\sqrt{a-1}+\sqrt{b-1}$. Godsil's bound \cite{Go} puts every root of $\mu_{I(H)}$ in
$[-\rho,\rho]$, and Proposition~\ref{prop:bridge} squares them.
\end{proof}

The coordinate hypothesis is basis-dependent as stated and is not meant to be. It says exactly
that the family commutes.

\begin{corollary}\label{cor:commuting}
Let $a,b\ge2$ and let $P_1,\dots,P_q$ be pairwise commuting rank-$b$ orthogonal projections on
$\mathbb R^n$ with $\sum_kP_k=aI$. Then $\operatorname{maxroot}\mu\le(\sqrt{a-1}+\sqrt{b-1})^2$, and
the constant is attained in the limit over this class.
\end{corollary}

The hypothesis $a,b\ge2$ is needed and not cosmetic. At $a=b=1$, $n=q=1$, $P_1=[1]$, one has
$\mu(x)=x-1$ and $\operatorname{maxroot}\mu=1$ against a right-hand side of $0$. The bound is a
statement about the $(a,b)$-biregular tree, and that tree is a single edge or a ray unless both
degrees are at least two.

\begin{proof}
Commuting orthogonal projections are simultaneously diagonalisable, and in a common eigenbasis
each $P_k$ is a $0/1$ diagonal matrix with exactly $b$ ones, that is a hyperedge $e_k$ on $n$
vertices; $\sum_kP_k=aI$ says every vertex lies in exactly $a$ of them. So the family is the
coordinate family of an $a$-regular $b$-uniform hypergraph, and Corollary~\ref{cor:xucoord} and
Proposition~\ref{prop:xusharp} apply. Both statements are invariant under a common orthogonal
change of basis, so nothing is lost in choosing one.
\end{proof}

So the inequality Xu conjectures is a theorem on the commuting locus, at every rank, and not only
at $b=2$ where it is Heilmann--Lieb. The question that matters for the conjecture is whether the
constant is the right one there, and it is.

\begin{proposition}\label{prop:xusharp}
Fix $a,b\ge2$. Over coordinate families with parameters $(a,b)$,
\[
\sup\ \frac{\operatorname{maxroot}\mu}{\bigl(\sqrt{a-1}+\sqrt{b-1}\bigr)^2}\ =\ 1 .
\]
\end{proposition}

\begin{proof}
At most $1$ by Corollary~\ref{cor:xucoord}. For the other direction, let $B_r$ be the ball of
radius $r$ in the $(a,b)$-biregular tree. If $\operatorname{girth}I(H)>2r$ then the $r$-ball of
$I(H)$ about any $a$-side vertex is acyclic and, by biregularity, isomorphic to $B_r$, so
$B_r\subseteq I(H)$ as a subgraph. The greatest root of a matching polynomial is monotone under
subgraphs \cite{HL,Go}, and $B_r$ is a tree, where the matching polynomial is the characteristic
polynomial; hence
$\operatorname{maxroot}\mu_{I(H)}\ge\lambda_{\max}(B_r)$, and squaring,
$\operatorname{maxroot}\mu\ge\lambda_{\max}(B_r)^2$. The balls exhaust the tree and
$\lambda_{\max}(B_r)\uparrow\rho$. Finally $(a,b)$-biregular bipartite graphs of arbitrarily large
girth exist: random $N$-lifts of any one of them preserve biregularity and have girth tending to
infinity in probability. Taking $r\to\infty$ gives the claim.
\end{proof}

The ladder is measured rather than assumed, and it also supplies an explicit floor for any graph
whose girth is known. Writing the ratio in the $\mu$ variable,
$\lambda_{\max}(B_r)^2/\rho^2$ is
\begin{center}
\begin{tabular}{lcccccc}
$(a,b)$ & $r=2$ & $4$ & $6$ & $8$ & $10$ & $12$\\
\hline
$(3,2)$ & $0.6863$ & $0.8579$ & $0.9167$ & $0.9443$ & $0.9596$ & $0.9692$\\
$(4,3)$ & $0.6061$ & $0.8082$ & $0.8854$ & $0.9234$ & $0.9450$ & $0.9585$\\
$(5,4)$ & $0.5744$ & $0.7898$ & $0.8747$ & $0.9168$ & $0.9396$ & \\
$(4,4)$ & $0.5833$ & $0.7951$ & $0.8777$ & $0.9186$ & $0.9418$ & $0.9505$\\
$(6,3)$ & $0.6004$ & $0.8031$ & $0.8826$ & $0.9219$ & $0.9442$ & $0.9489$\\
\hline
\end{tabular}
\end{center}
increasing in $r$ at every pair, $b\ge3$ included; the missing entry is where the ball passed the
$400{,}000$ vertex cap, and a truncated ball is a proper subtree of the true one, hence still a
subgraph of any graph of girth $>2r$, so the floor it gives remains valid. On named incidence
graphs, with girth recomputed by breadth-first search rather than quoted, the realized ratios are
$0.8995$ for Heawood, $0.9171$ for Pappus and $0.9453$ for Tutte--Coxeter at $(3,3)$; $0.9098$ for
$\mathrm{AG}(2,3)$ at $(4,3)$; and $0.9097$ for $\mathrm{PG}(2,3)$ at $(4,4)$. Every one exceeds
its ball floor, which is the containment argument of Proposition~\ref{prop:xusharp} checked
against the object it is about (\texttt{code/xu\_sharp.py}).

Two things follow, and the second is the reason for the subsection.

First, the plateaus found by the adversarial search of \texttt{code/adversarial.py} and
\texttt{code/rl\_push.py} --- $0.739$ at $(4,3)$ and $0.726$ at $(5,4)$, against $0.939$ at
$(3,2)$, none of them a violation --- are an artefact of that search, not a property of the
class. It sampled generic subspaces, and generic subspaces behave like short
girth. The coordinate construction reaches $0.9098$ at $(4,3)$ on twenty-one vertices and is
bounded below by $0.9585$ at girth $26$, well past where the search stalled. There is no evidence of a smaller true constant at
$b\ge3$, and the reading that suggested one is closed.

Second, the constant is now fixed. Any proof of Xu's conjecture in general must produce exactly
$(\sqrt{a-1}+\sqrt{b-1})^2$; any attempted improvement of it is false, by
Proposition~\ref{prop:xusharp}, before it is attempted. Equally, the whole content of the
conjecture lies off the commuting locus: the bridge of Proposition~\ref{prop:bridge} fails for
non-coordinate rank-$b$ projections, which is checked as a control rather than asserted, and it is
that failure which leaves the general statement open.

Whether that leaves the general statement hard or merely unreduced turns on one further
question, which is measurable: is the commuting locus \emph{extremal}? Deform a commuting family
by rotating each block independently, $U_e=\exp(tX_e)$ with $X_e$ skew, and return to the tight
projection class by alternating projections onto $\{$rank-$b$ projections$\}$ and
$\{\sum_kA_k=aI\}$. Along such a path the ratio falls at every hypergraph tried; an ascent
started at the commuting point accepted none of $400$ steps at any of four starting points; and,
which is the part that bears on a global statement, no random start in the tight projection class
beat its commuting counterpart at the same $(n,q,a,b)$, over fifteen starts each, the best random
family falling short by between $0.0038$ and $0.0107$ (\texttt{code/offcoord.py}). On that
evidence Conjecture~1.4 would follow from Corollary~\ref{cor:commuting} together with the
statement that the commuting locus is extremal, which replaces a conjecture about all tight
rank-$b$ families by a monotonicity statement over a theorem.

Proposition~\ref{prop:singular} says which directions are excluded; the same computation says
which are admitted, and the answer is an explicit cone.

Call $i$ and $j$ \emph{twins} when they lie in exactly the same hyperedges, and the family
\emph{twin-free} when it has no twin pair. Twin-freeness is checkable in $O(n^2q)$ and holds for
every family used below; the Fano family, $C_4$ and the triples of $K_4$ are all twin-free.

\begin{proposition}\label{prop:cone}
Let $D$ lie in the kernel of the linearised constraints. A second-order correction $X$ exists if and
only if
\[
Q_j(D):=\sum_k\sigma_k(j)\,(D_k^2)_{jj}=0\qquad\text{for every vertex }j,
\]
with $\sigma_k(j)$ as in Proposition~\ref{prop:singular}, \emph{together with}
\[
\sum_k\sigma_k(i,j)\,(D_k^2)_{ij}=0\qquad\text{for every twin pair }\{i,j\},
\qquad \sigma_k(i,j)=1-(P_k)_{ii}-(P_k)_{jj}.
\]
For a twin-free family the second family of conditions is empty, and the directions admitting a
second-order correction are then the kernel of the linearisation cut by the $n$ quadrics $Q_j$.
\end{proposition}

\begin{proof}
The order-two equation $X_k-P_kX_k-X_kP_k=D_k^2$ has $(i,j)$ entry
$(X_k)_{ij}\bigl(1-(P_k)_{ii}-(P_k)_{jj}\bigr)=(D_k^2)_{ij}$, so the coefficient is $-1$ when both
$i,j$ lie in $e_k$, $+1$ when neither does, and $0$ when $e_k$ splits the pair. Where the coefficient
is nonzero $(X_k)_{ij}$ is forced; where it vanishes the entry is free and the equation demands
$(D_k^2)_{ij}=0$, which holds automatically since $D_k$ is supported on split pairs and $D_k^2$
therefore is not. Imposing $\sum_kX_k=0$ entry by entry, a pair is solvable exactly when some $e_k$
splits it, leaving a free entry to absorb the forced sum, or else the forced values already cancel.
No $e_k$ splits a diagonal pair, which gives the conditions $Q_j(D)=0$; and no $e_k$ splits an
off-diagonal pair exactly when its two vertices are twins, which gives the second family.
\end{proof}

The entrywise criterion is machine-checked as \texttt{ConeFix.solvable\_iff}, together with
\texttt{coef\_eq\_zero\_iff} for the combinatorial reading of the coefficient and
\texttt{no\_free\_entry\_iff\_twin} for the identification of the obstructed pairs as the twin
pairs. The twin conditions are not vacuous. On the $3$-uniform $2$-regular family with vertices
$\{0,\dots,5\}$ and hyperedges $\{0,1,2\},\{0,1,3\},\{2,4,5\},\{3,4,5\}$, where $0$ and $1$ are
twins, the direction $D_1=E_{04}+E_{40}+E_{14}+E_{41}$, $D_2=0$, $D_3=-(E_{04}+E_{40})$,
$D_4=-(E_{14}+E_{41})$ lies in the kernel and satisfies $Q_j(D)=0$ at all six vertices, while the
forced values at the twin entry $(0,1)$ sum to $-1$: no second-order correction exists. Every step of
that witness is machine-checked in \texttt{RamaLean/ConeWitness.lean}
(\texttt{twins}, \texttt{cross\_supported}, \texttt{sum\_D\_eq\_zero}, \texttt{Q\_eq\_zero},
\texttt{twin\_obstruction}), so the necessity of the twin conditions rests on no arithmetic done
outside the kernel.

\begin{proof}[Proof of the twin-free case, restated]
Necessity is the computation in Proposition~\ref{prop:singular}: the diagonal of each $X_k$ is
forced and summing it over $k$ gives $Q_j(D)$, which $\sum_kX_k=0$ requires to vanish. For
sufficiency, every off-diagonal pair is split by some $e_k$ and the free entry there absorbs the
forced sum, so only the diagonal conditions remain.
\end{proof}

The second-order behaviour of the top root over that cone can now be examined, and it is better
than the cone alone would suggest. At a critical point the second derivative of $\lambda_{\max}$
along a curve in $V$ depends only on the curve's velocity, so it defines a quadratic form $\Lambda$
on the admissible directions; and since the $2$-jet $P_k+\varepsilon D_k+\varepsilon^2X_k$ agrees
with a curve on $V$ to $O(\varepsilon^3)$, $\Lambda$ may be read off that jet without constructing
the curve. Computed this way on the Fano family, by fitting the coefficients of $\mu$ as
polynomials in $\varepsilon$ rather than differencing them, $\Lambda$ is negative semidefinite on
the whole linearised kernel, with $42$ negative eigenvalues, $21$ zero and none positive, and its
null space is exactly the tangent space to the orthogonal conjugation orbit
(\texttt{code/hessian.py}). The value on a cross direction agrees to eight digits with the exact
rational $-C/4$ of Proposition~\ref{prop:rigid}'s computation.

Two things follow. First, no argument restricted to the cone is needed for the second-order test:
$\Lambda\preceq0$ already on the kernel, so it is negative on the cone a fortiori, strictly away
from the conjugation directions along which $\mu$ does not change at all. Second, the commuting
family is a strict local maximum of $\lambda_{\max}$ to second order, transverse to its own orbit.

The signature is certified rather than observed. Each entry of the form is an exact integer
polynomial in $y$: the jet has integer entries when $D$ does, which the basis arranges, so the
coefficients of $\mu$ are integer polynomials in $\varepsilon$ and the $\varepsilon^2$ part is
extracted by truncated series arithmetic over $\mathbb Z$. The vanishing on the orbit needs no
arithmetic at all, since $\mu$ is invariant under simultaneous orthogonal conjugation and
$\lambda_{\max}$ is therefore constant along $\exp(\varepsilon\Omega)P_k\exp(-\varepsilon\Omega)$;
the computation returns exactly $0$ there. On a complement of the orbit the restricted form has
eigenvalues in $[-0.1269,-0.0218]$ and its negative admits a Cholesky factorisation with smallest
pivot $3.7\times10^{-2}$, against an evaluation error of order $10^{-55}$ at the working precision
(\texttt{code/curvform.py}). The margin exceeds the error by some $10^{53}$, so the definiteness
is a certificate and not a reading.

The extra null directions that appear at $b=2$ are not a failure of maximality. Beyond the
conjugation orbit the form has $2$ further null directions for $C_4$, $6$ for $K_4$ and $9$ for
$C_6$, all exact, and none at $(2,3)$, $(3,3)$ or the family of triples in $K_4$; so the trigger is
$b=2$ and not $a=b=2$, $K_4$ and $K_{3,3}$ being $(3,2)$. Along those directions
$\lambda_{\max}$ still decreases, at fourth order rather than second, with the same sign for both
signs of the parameter: the fourth-order coefficients measured are $-0.0038$ and $-0.0005$ for
$C_4$ and between $-0.60$ and $-0.14$ for $K_4$. At $b=2$ the second-order test is blind and the
fourth order does the work.

It is not a statement about all commuting families, and the qualification is not idle. The same
computation on four further families gives no positive eigenvalue anywhere, so $\Lambda\preceq0$
throughout; but the strictness transverse to the orbit fails at $(a,b)=(2,2)$, where for $C_6$ the
null space has dimension $24$ against an orbit of rank $15$, nine directions carrying no curvature
without being conjugations. It holds at $(2,3)$ with smallest pivot $8.7\times10^{-2}$ and at
$(3,3)$, and is vacuous for the family of triples in $K_4$, whose whole kernel is the orbit. The
$(2,2)$ case is the one whose universal cover is a line, with no spectral gaps at all, so its
degeneracy is not surprising; but on the evidence the correct statement is $\Lambda\preceq0$ in
general, with strictness established only at the families where it was certified. The excess is
triggered by $b=2$ and not by $a=b=2$: $K_4$ and $K_{3,3}$ are $(3,2)$ and carry excess $6$ and $9$,
as the table above records.

One order beyond the obstruction, nothing further is imposed, and for a reason of supports. The
order-three equation reads $Y_k-P_kY_k-Y_kP_k=D_kX_k+X_kD_k$. Its off-diagonal blocks give the
condition $(D_kX_k+X_kD_k)_{12}=0$, which holds identically: with $X_{11}=-(D^2)_{11}$ and
$X_{22}=(D^2)_{22}$ from order two, that block is
$D_{12}D_{21}D_{12}-D_{12}D_{21}D_{12}=0$. Its diagonal entries vanish too. The diagonal is where the
content is. No $e_k$ splits a diagonal pair, so $(Y_k)_{jj}$ is forced and $\sum_kY_k=0$ requires
\[
R_j:=\sum_k\sigma_k(j)\,(D_kX_k+X_kD_k)_{jj}=0\qquad\text{for every }j .
\]
$R$ is not identically zero: at the Fano family a second-order correction satisfying $\sum_kX_k=0$
gives $\max_j|R_j|=19.7$. What is identically zero is its \emph{sum}, and that is a two-line
theorem holding per block, with no hypothesis on $X$ at all.

\begin{lemma}[order-three trace law]\label{lem:o3trace}
If $P_kD_k+D_kP_k=D_k$ then $\operatorname{tr}(D_kX_k+X_kD_k)=2\operatorname{tr}\bigl(P_k(D_kX_k+X_kD_k)\bigr)$
for every $X_k$, that is $\sum_j\sigma_k(j)(D_kX_k+X_kD_k)_{jj}=0$. Summing over $k$, $\sum_jR_j=0$.
\end{lemma}

\begin{proof}
$\operatorname{tr}(P_k(D_kX_k+X_kD_k))=\operatorname{tr}(P_kD_kX_k)+\operatorname{tr}(P_kX_kD_k)
=\operatorname{tr}(P_kD_kX_k)+\operatorname{tr}(D_kP_kX_k)=\operatorname{tr}((P_kD_k+D_kP_k)X_k)
=\operatorname{tr}(D_kX_k)$ by cyclicity and the order-one equation, while
$\operatorname{tr}(D_kX_k+X_kD_k)=2\operatorname{tr}(D_kX_k)$.
\end{proof}

Machine-checked as \texttt{OrderThreeLaw.trace\_order\_three}. It needs neither idempotency, nor
tightness, nor any support hypothesis, and it replaces a complementary-supports argument given in an
earlier version of this note whose hypothesis on $X$ is satisfied by no second-order correction: at
the Fano family the canonical $X$ with vanishing split-pair entries fails $\sum_kX_k=0$ at all $21$
of $21$ off-diagonal pairs, those entries being the only free ones and hence the only ones that can
absorb the forced same-side sums.

The second-order correction is not unique, its split-pair entries being free, and $R$ depends on the
choice. So order three is unobstructed exactly when some admissible $X$ makes $R=0$. Lemma
\ref{lem:o3trace} puts $R$ in the hyperplane $\sum_jR_j=0$ for every choice, and the image of the
free entries is the same object that governs the second order.

\begin{proposition}\label{prop:o3rank}
Write $u$ for the free entries of $X$, which satisfy $\sum_{k\in K(i,j)}u_{k,i,j}=0$ pair by pair.
The map $u\mapsto R$ is linear with no constant term and has rank $n-c(G_D)$, with $G_D$ the graph of
Proposition~\ref{prop:cone}'s companion: $i\sim j$ when $(\sigma_k(j)D_k(i,j))_{k\in K(i,j)}$ is not
constant in $k$. Consequently, if $G_D$ is connected then order three is unobstructed.
\end{proposition}

\begin{proof}
$R$ is linear in $u$ because $D_k$ is supported exactly on the split pairs, which are exactly the
free entries. Using $\sigma_k(i)=-\sigma_k(j)$ on a split pair,
\[
\sum_jw_jR_j=2\sum_{i<j}(w_j-w_i)\!\!\sum_{k\in K(i,j)}\!\!\sigma_k(j)D_k(i,j)u_{k,i,j},
\]
so $w$ annihilates the image exactly when, pair by pair, either $w_i=w_j$ or the linear functional
$u\mapsto\sum_k\sigma_k(j)D_k(i,j)u_k$ vanishes on $\{\sum u=0\}$, which happens exactly when its
coefficients are constant in $k$. That is the definition of $G_D$, and the annihilator is therefore
the functions constant on its components, of dimension $c(G_D)$. If $G_D$ is connected the rank is
$n-1$, which is the dimension of the hyperplane $\sum_jR_j=0$ that Lemma~\ref{lem:o3trace} confines
$R$ to, so the image is that hyperplane and $R=0$ is attained.
\end{proof}

Every step is machine-checked: \texttt{OrderThreeLaw.trace\_order\_three} for the lemma,
\texttt{SpanRank.linear\_vanishes\_iff\_const} for the annihilator condition,
\texttt{CokernelRank.const\_of\_adj\_eq} for the passage from edges to components, and
\texttt{OrderFourSolvable.mem\_of\_sum\_zero} for the conclusion. The rank formula is verified
against the computed rank at five families and twenty cone directions, and order three is solvable at
every one of them, including $C_4$, whose $G_D$ has two components so that the proposition does not
apply and solvability there is measured (\texttt{code/orderthree.py}).

Order four is then satisfied as well, and the freedom in $Y$ is what absorbs it. The condition is
$\sum_k\sigma_k(j)\bigl[2(D_kY_k)_{jj}+(X_k^2)_{jj}\bigr]=0$ at each vertex, linear in $Y$, so it
holds exactly when the $X^2$ vector lies in the image of $L_D:Y\mapsto
\bigl(2\sum_k\sigma_k(j)(D_kY_k)_{jj}\bigr)_j$. That image always lies in the trace-zero hyperplane, and for the same
kind of reason: writing the corners of $D_k$ and $Y_k$ as $M$ and $N$, the two halves of
$\operatorname{tr}\bigl((I-2P_k)D_kY_k\bigr)$ are $\operatorname{tr}(MN^{\mathsf T})$ and
$\operatorname{tr}(M^{\mathsf T}N)$, which are equal, so the sum over $j$ of the image vanishes
identically (\texttt{OrderFour.image\_in\_trace\_zero}). That it fills that hyperplane has a criterion. A vector $w$
annihilates the image exactly when $(w_i-w_j)\bigl(\sigma_k(i)D_k(i,j)-\sigma_l(i)D_l(i,j)\bigr)=0$
for every pair $i,j$ and every pair of blocks $k,l$ separating them, so $w_i=w_j$ as soon as two
separating blocks disagree at $(i,j)$. Joining $i$ to $j$ whenever they do, $w$ is constant on each
component and
\[
\operatorname{rank}L_D=n-c(G_D),
\]
with $c$ the number of components. The cokernel is therefore spanned by $(1,\dots,1)$, and order
four is solvable, exactly when $G_D$ is connected. Checked against the computed rank on $32$
directions across two families, agreeing in every case, the cross basis directions being precisely
those where $G_D$ is empty and the rank collapses to $1$; there the obstruction vanishes outright
and order four is solvable regardless. The scalar step and the passage from edges to components are
machine-checked as \texttt{CokernelRank.sep\_forces\_eq} and
\texttt{CokernelRank.const\_of\_adj\_eq}. And the
$X^2$ vector lies in it for a reason. Summing over the vertices,
\[
\sum_j\sum_k\sigma_k(j)(X_k^2)_{jj}
=\sum_k\Bigl[\operatorname{tr}\bigl((D_k^2)_{22}^2\bigr)-\operatorname{tr}\bigl((D_k^2)_{11}^2\bigr)\Bigr],
\]
the contributions of the free block $X_{12}$ cancelling through
$\operatorname{tr}(X_{21}X_{12})=\operatorname{tr}(X_{12}X_{21})$. Since $D_k$ is off-diagonal for
the splitting, writing its corner as $M$ gives $(D_k^2)_{11}=MM^{\mathsf T}$ and
$(D_k^2)_{22}=M^{\mathsf T}M$, whose squares have equal trace by cyclicity, so every term vanishes.
The identity is machine-checked as \texttt{OrderFour.order\_four\_sum\_zero}, and the residuals of
the order-four solve are $10^{-18}$ on the directions tested (\texttt{code/tangentcone.py}). Those
two facts, that the obstruction sums to zero and that the cokernel is the constants, do not by
themselves put the obstruction in the image, and the step that does is
\texttt{OrderFourSolvable.order\_four\_solvable}: in finite dimension a subspace whose orthogonal
complement is the line of constants is the hyperplane of vectors summing to zero. It had been left
unstated between two files that each proved half of it.

Orders two, three and four are therefore all satisfiable on the cone. The passage from every finite
order to an actual curve was the remaining gap, and it can be closed rather than climbed, because
the variety is in the scope of a classical criterion.

Beyond the second order, the extremality itself has been searched for directly. At four
commuting families the commuting point is the largest value found over between $32$ and $40$
random tight projection families at the same $(p,q,a,b)$, by margins of $0.0004$ for $C_4$,
$0.016$ for $C_6$, $0.046$ for $K_{3,3}$ and $0.048$ for the cube, every value staying below the
band edge (\texttt{code/extremal.py}). That is the statement of \S\ref{sec:rank} seen from the
other side: within the projection class the commuting families are the largest, while across the
rank-$n$ interpolation the projections are the smallest. Neither observation is a proof, and the
conjunction of the two is what Conjecture~1.4 would need.

\begin{proposition}[the curve, by 2-regularity]\label{prop:tworeg}
Let $D$ lie in the kernel of the linearisation with $Q(D)=0$, and suppose the linear map
$E\mapsto\mathrm dQ(D)[E]$ carries that kernel onto the span of $Q_1,\dots,Q_n$. Then there is a
curve $t\mapsto A(t)$ in the tight projection variety with $A(0)=A_0$ and $A(t)=A_0+tD+O(t^2)$. In
particular $D$ lies in the tangent cone.
\end{proposition}

\begin{proof}
Work upstairs, on frames $U_k\in\mathbb R^{n\times b}$ with $U_k^{\mathsf T}U_k=I_b$ and
$A_k=U_kU_k^{\mathsf T}$, so that idempotency is free and the only equation is
$\Psi(U)=\sum_kU_kU_k^{\mathsf T}-aI=0$. The Stiefel constraint makes the $e_k$-block of a
first-order frame variation antisymmetric, and a diagonal entry of $\mathrm d\Psi$ is a diagonal
entry of that block, so $\mathrm d\Psi$ has vanishing diagonal; and it is onto the zero-diagonal
symmetric matrices, of rank $21$ against $28$ at the Fano family. Split $\Psi$ accordingly. Its
off-diagonal part is a submersion, so $W=\{\Psi_{\mathrm{off}}=0\}$ is a smooth manifold near the
point, and on $W$ the remaining map $F=\Psi_{\mathrm{diag}}$ satisfies $F=0$ and $\mathrm dF=0$
there, with Hessian $Q$.

Write $F=Q+C+\cdots$ with $Q(x)=B(x,x)$ quadratic and $C$ cubic. Since $Q(D)=0$, the polynomial
$t\mapsto F(tD+t^2w)$ vanishes to order three, so $F(tD+t^2w)=t^3G(t,w)$ with $G$ polynomial and
$G(0,w)=2B(D,w)+C(D,D,D)$, affine in $w$ with surjective linear part by hypothesis. Pick $w_0$ with
$G(0,w_0)=0$ and apply the implicit function theorem in $w$ at $(0,w_0)$; the resulting
$x(t)=tD+t^2w(t)$ satisfies $F(x(t))=0$ and $x(t)/t\to D$.
\end{proof}

This is Avakov's 2-regularity condition \cite{Avakov,Arutyunov} in the fully degenerate case, and
no novelty is claimed for it. What is new here is only that the tight projection variety is in its
scope, which is what the frame reduction supplies.

The span in the hypothesis is not all of $\mathbb R^n$, and the reason is an identity that holds for
each block separately,
\[
\sum_j\sigma_k(j)(D_k^2)_{jj}=0 ,
\]
since $(D_k^2)_{jj}=\sum_lD_k(j,l)^2$ by symmetry, so that exchanging the two indices in the double
sum pairs every surviving term with one of opposite sign, $D_k$ being supported on pairs with
exactly one endpoint in $e_k$. So $Q$ lands in the trace-zero diagonals, and the span can be smaller
still: it is $n-1$ at $C_6$, $K_{3,3}$ and the Fano family, and $2$ at $C_4$, which carries the extra
pair $Q_0+Q_2=Q_1+Q_3=0$, one for each side of its bipartition. The identity is machine-checked as
\texttt{ArcTangent.sum\_sgn\_diag\_sq\_eq\_zero}, and the passage from a curve to membership of the
tangent cone as \texttt{ArcTangent.mem\_tangentCone\_of\_arc}; the implicit function theorem step is
not formalised.

The hypothesis is a rank condition on $D$ alone, and it holds where it has been checked: at every
cone direction tested at $C_4$, $K_{3,3}$ and the Fano family, four directions each, the rank of
$\mathrm dQ(D)$ on the kernel equals the span. The curve has also been exhibited rather than
inferred, by continuation upstairs: the residual stays at $10^{-14}$ or below and
$\|(A(t)-A_0)/t-D\|$ falls with exponent $0.96$ to $1.00$ over $t\in[0.01,0.05]$, while a direction
with $Q(D)\neq0$ has the same quantity flat at exponent $0$ (\texttt{code/arc.py}).

Both quantities in that hypothesis are combinatorial, and saying so answers the question of when it
holds. Write $K(i,j)$ for the blocks separating $i$ from $j$. A kernel direction has $D_k$ supported
on separated pairs and tightness reads $\sum_{k\in K(i,j)}D_k(i,j)=0$ pair by pair, with distinct
pairs independent; and since $(D_k^2)_{jj}=\sum_lD_k(j,l)^2$ and $\sigma_k(i)=-\sigma_k(j)$ for
separating $k$,
\[
\sum_jw_jQ_j(D)=\sum_{i<j}(w_j-w_i)\!\!\sum_{k\in K(i,j)}\!\!\sigma_k(j)D_k(i,j)^2,\qquad
\sum_jw_j\,\mathrm dQ_j(D)[E]=2\sum_{i<j}(w_j-w_i)\!\!\sum_{k\in K(i,j)}\!\!\sigma_k(j)D_k(i,j)E_k(i,j).
\]
So in each case $w$ annihilates exactly when, pair by pair, either $w_i=w_j$ or a form vanishes
identically on the hyperplane $\{\sum t=0\}$. A linear functional vanishes there exactly when its
coefficients are constant; a diagonal quadratic form $\sum\varepsilon_kt_k^2$ vanishes there exactly
when $\varepsilon_a+\varepsilon_b=0$ for every pair, which for signs $\pm1$ means
$|K(i,j)|\le1$, or $|K(i,j)|=2$ with the two separating blocks on opposite sides. Writing $G_Q$ for
the graph on the vertices joining $i$ to $j$ when the quadratic form does \emph{not} vanish, and
$G_D$ for the graph joining them when $(\sigma_k(j)D_k(i,j))_{k\in K(i,j)}$ is not constant,
\[
\operatorname{span}\{Q_1,\dots,Q_n\}=n-c(G_Q),\qquad
\operatorname{rank}\mathrm dQ(D)=n-c(G_D),
\]
with $c$ the component count. $G_Q$ depends on the hypergraph alone, $G_D$ on the direction; $G_D$ is
always a subgraph of $G_Q$, so the rank never exceeds the span, and 2-regularity at $D$ says exactly
that the two graphs have the same components. Both formulas agree with the computed values at seven
families (\texttt{code/spanrank.py}), and the scalar steps are machine-checked in
\texttt{RamaLean/SpanRank.lean}.

On a tight family the criterion collapses further, and the collapse says exactly when the span drops.
Writing $\lambda_{ij}$ for the number of blocks containing both $i$ and $j$,
\[
|K(i,j)|=\deg(i)+\deg(j)-2\lambda_{ij}=2(a-\lambda_{ij}),
\]
every degree being $a$. So $|K(i,j)|$ is even and the two sides of a separated pair carry
$a-\lambda_{ij}$ blocks each; the case $|K|=2$ with both blocks on the same side cannot arise, and
$|K|=2$ always means the cross configuration. The whole criterion becomes a codegree condition,
\[
G_Q:\quad i\sim j\iff\lambda_{ij}\le a-2,
\]
so $G_Q$ is the complement of the graph joining pairs of codegree at least $a-1$, and
\begin{center}
the span drops below $n-1$ $\iff$ the codegree-at-least-$(a-1)$ graph is a complete join,
\end{center}
a graph being a join exactly when its complement is disconnected. Over a sweep of tight families
this happens four times (\texttt{code/spanrank.py}): at $C_3$, whose codegree graph is $K_3$; at
$C_4$, whose codegree graph is $C_4=K_{2,2}$, giving the complement $2K_2$ and $c(G_Q)=2$, which is
the exception recorded above; at the four triples of $K_4$, where every pair lies in two of them, so
the codegree graph is $K_4$, the quadrics vanish identically and the tangent cone is the whole
kernel; and at the $2$-regular $3$-uniform family on six vertices, whose codegree graph is the
cocktail party graph $K_{2,2,2}$. For $a=b=2$ the codegree graph is the cycle itself, and among
cycles only $C_3$ and $C_4$ are joins, which is why every longer cycle has span exactly $n-1$.

It also settles, in the negative, whether every cone direction is 2-regular. A cross basis direction
is supported on a single pair with two separating blocks of opposite sign; the vanishing criterion
above is exactly why it satisfies $Q=0$, so it lies on the cone, and the constancy criterion makes
$\operatorname{rank}\mathrm dQ(D)$ equal to $0$ when $|K(i,j)|=2$ and to $1$ otherwise, against a span
of at least $2$. Checked at seven families, where the number of such directions runs from $12$ at
$C_4$ to $324$ at $\mathrm{AG}(2,3)$. Those directions are not lost: Proposition~\ref{prop:rigid}'s
rotation is an explicit curve with exactly that velocity. What is lost is a single uniform route,
and the two routes together cover every direction tested.

The earlier evidence, which measured distances rather than curves, is consistent and is recorded
because it covers directions the rank condition was not tested on. Over the $42$ cross basis
directions of the Fano family and
over directions obtained by projecting random kernel elements onto $\{Q=0\}$, the nearest point of
the variety to $A_0+\varepsilon D$ lies at distance $O(\varepsilon^2)$, and at $O(\varepsilon)$ for
every direction with $Q(D)\neq0$ (\texttt{code/tangentcone.py}); the general relation is
$\operatorname{dist}(A_0+\varepsilon D,V)=\varepsilon\operatorname{dist}(D,T)+O(\varepsilon^2)$,
which is what the two exponents measure. The consequence worth recording is that second-order
optimisation over the commuting locus is a constrained problem with $n$ quadratic constraints
rather than an unavailable one.

The deformation is also far more rigid than the tangent cone alone suggests.

\begin{proposition}\label{prop:rigid}
Along the rotation of Proposition~\ref{prop:cone}'s cross configuration, the coefficients of
$y^{n},y^{n-1},y^{n-2},y^{n-3}$ in $\mu$ are independent of $\theta$, so the $\theta$-dependent
part of $\mu$ has degree at most $n-4$ in $y$.
\end{proposition}

\begin{proof}
Only subsets containing exactly one of $e,f$ carry $\theta$, and they pair as $S\cup\{e\}$ with
$S\cup\{f\}$. Write $N_e=P_{e\setminus v}+M_S$ and $N_f=P_{f\setminus w}+M_S$, both diagonal since
every block other than $e,f$ is a coordinate projection, and put $u=\cos\theta\,\varepsilon_v+
\sin\theta\,\varepsilon_w$, $u^{\perp}=-\sin\theta\,\varepsilon_v+\cos\theta\,\varepsilon_w$, so
that $A_e+M_S=N_e+P_u$ and $A_f+M_S=N_f+P_{u^\perp}$. From
$\det(xI+N+P_u)=\det(xI+N)+u^{\mathsf T}\operatorname{adj}(xI+N)u$ and the diagonality of $N$,
whose adjugate has $(i,i)$ entry $\prod_{l\neq i}(x+d_l)$,
\[
\det(xI+A_e+M_S)+\det(xI+A_f+M_S)=\bigl[\theta\text{-free}\bigr]
+\sin^2\theta\,\Delta(x),
\]
\[
\Delta(x)=(d_v-d_w)\Bigl[\prod_{l\neq v,w}(x+d_l)-\prod_{l\neq v,w}(x+d'_l)\Bigr],
\]
where $d=\operatorname{diag}N_e$, $d'=\operatorname{diag}N_f$, using $d_v=d'_v$ and $d_w=d'_w$,
both of which hold because $v\notin f$ and $w\notin e$. The two products are monic of degree
$n-2$, and $\sum_ld_l=\operatorname{tr}N_e=(b-1)+\operatorname{tr}M_S=\operatorname{tr}N_f=
\sum_ld'_l$, so they agree in their coefficients of $x^{n-2}$ and $x^{n-3}$. Their difference
therefore has degree at most $n-4$, and so does $\Delta$.
\end{proof}

The bound is attained: the $\theta^2$ part of $\mu$ has degree exactly $n-4$ for the Fano family,
where it is $2y(y-1)(y-4)$, and for $\mathrm{PG}(2,3)$, where it has degree $9$
(\texttt{code/curvature.py}). The step of the proof carrying the exponent, that two monic
polynomials with equal root sums differ in degree at most $n-4$, is machine-checked as
\texttt{CoefficientRigidity.prod\_diff\_natDegree\_le}.

The second-order behaviour can be computed exactly rather than sampled, along curves that stay on
the variety identically. Take hyperedges $e,f$ and coordinates $v\in e\setminus f$,
$w\in f\setminus e$, let $R(\theta)$ rotate the plane $\operatorname{span}\{\varepsilon_v,
\varepsilon_w\}$, and apply it to $P_e$ and $P_f$ alone. Since $P_v+P_w$ is the identity on that
plane, $R(P_e+P_f)R^{\mathsf T}=P_e+P_f$ identically in $\theta$, so tightness is preserved
exactly and not merely to leading order, while each rotated block stays an orthogonal projection
of rank $b$ and the family genuinely decommutes. For the Fano family the mixed characteristic
polynomial is then computed in exact rational arithmetic, its odd part vanishes as the symmetry
$\theta\mapsto-\theta$ requires, and
\[
\lambda_{\max}(\theta)=y_0-C\theta^2+O(\theta^4),\qquad
y_0=7.19605818500856\ldots,\quad C=0.03717851401295\ldots,
\]
with $y_0$ a root of $y^7-21y^6+168y^5-644y^4+1218y^3-1050y^2+336y-24$ and $C$ algebraic of the
same degree. All $168$ curves of this shape give the same $C$, as the $2$-transitivity of the Fano
plane requires.

Running the same computation on two further commuting families shows what $C$ depends on, and it
is not what one would guess.
\begin{center}
\begin{tabular}{lccccc}
family & $(a,b)$ & $n$ & $y_0$ & $C$, $|e\cap f|=1$ & $C$, $|e\cap f|=0$\\
\hline
Fano & $(3,3)$ & $7$ & $7.196058185$ & $0.037178514013$ & ---\\
$\mathrm{AG}(2,3)$ & $(4,3)$ & $9$ & $9.005835502$ & $0.021354509344$ & $0.041740312350$\\
$\mathrm{PG}(2,3)$ & $(4,4)$ & $13$ & $10.916978461$ & $0.022914015843$ & ---\\
\hline
\end{tabular}
\end{center}
There is no closed form $C(a,b)$, and not for want of data: $\mathrm{AG}(2,3)$ carries two values
at the single parameter pair $(4,3)$, according to whether the two hyperedges rotated are disjoint
or meet. In the two projective planes every pair of lines meets, which is why each of them shows
one value; the affine plane has parallel classes and shows both. So the curvature is a function of
the pair of hyperedges, not of $(a,b)$, and the quantity a scaling law would name does not exist.
What is uniform is arithmetic rather than metric: in all three families $\mu_0$ is irreducible of
degree $n$ and $\deg C=\deg y_0=n$, so $C$ generates the same field as the top root
(\texttt{code/curvature.py}).

That is one diagonal entry of a second-order form, not definiteness, and the natural next step is
unavailable for a reason worth recording. Linearising idempotency and tightness at the commuting
point gives a space of dimension $63$, spanned by differences $m^k_{ij}-m^l_{ij}$ over hyperedges
containing exactly one of $i,j$. The variety does not fill it. Locating the nearest point of the
variety to $A_0+\varepsilon D$, the $42$ directions whose two hyperedges lie on opposite sides of
the pair --- exactly the tangents of the curves above --- sit at distance $O(\varepsilon^2)$ and
are tangent, while the remaining $21$ sit at distance $O(\varepsilon)$ and are not
(\texttt{code/hessian.py}). The commuting point is therefore a \emph{singular} point of the
variety. Its tangent cone is moreover not a linear subspace: $Q$ is quadratic, so it does not
vanish on sums, and four of the ninety-one pairs of cross basis directions tested have
$Q(D_1+D_2)\neq0$, reaching $2$. A Hessian on the linearised space is therefore the wrong
object. Local maximality therefore remains open, and what it needs is a second-order theory
on that cone.

The same is true at $\mathrm{AG}(2,3)$, where the linearised space has dimension $180$ against a
conjugation orbit of $36$: the $131$ cross directions give $\mathrm{dist}/\varepsilon^2$ equal to
$1.963$ and $1.998$ at $\varepsilon=0.08$ and $0.02$, and the $49$ same-group directions give
$\mathrm{dist}/\varepsilon$ equal to $1.4128$ and $1.4107$. The exponents agree with the Fano
family and only the constant differs.

Neither measurement is needed, as it turns out: the obstruction is algebraic and two blocks wide.

\begin{proposition}\label{prop:singular}
Let $\{P_k\}$ be a commuting tight family of rank-$b$ projections with hyperedges $\{e_k\}$, and
suppose there are $e,f$ and vertices $v,w$ with $w\in e\cap f$ and $v\notin e\cup f$. Then the
direction $D$ with $D_e=E_{vw}+E_{wv}$, $D_f=-(E_{vw}+E_{wv})$ and $D_k=0$ otherwise lies in the
kernel of the linearised constraints but is not tangent to the variety. Consequently the family is
a singular point of it, and the tangent cone is not a linear subspace.
\end{proposition}

\begin{proof}
Write $A_k=P_k+\varepsilon D_k+\varepsilon^2X_k+O(\varepsilon^3)$. Idempotency at order two is
$X_k-P_kX_k-X_kP_k=D_k^2$, whose $(i,j)$ entry for a coordinate projection is
$X_{ij}\bigl(1-(P_k)_{ii}-(P_k)_{jj}\bigr)$. At a diagonal entry $j$ the coefficient is
$1-2(P_k)_{jj}$, which is $-1$ when $j\in e_k$ and $+1$ when $j\notin e_k$ and is never zero, so
\[
(X_k)_{jj}=\sigma_k(j)\,(D_k^2)_{jj},\qquad
\sigma_k(j)=\begin{cases}-1,& j\in e_k,\\ +1,& j\notin e_k.\end{cases}
\]
The diagonal of $X_k$ is therefore determined; the freedom the equation leaves lies in the
off-diagonal blocks, which never meet a diagonal entry. Now $D_e^2=D_f^2=E_{vv}+E_{ww}$, the sign
squaring away, and $D_k^2=0$ for every other $k$. Since $v$ lies in neither $e$ nor $f$,
$\sigma_e(v)=\sigma_f(v)=+1$, so
\[
\sum_k(X_k)_{vv}=1+1+0+\cdots+0=2,
\]
contradicting $\sum_kX_k=0$; at $w$, which lies in both, the same computation gives $-2$. Hence no
$X$ exists, so no curve on the variety has velocity $D$. That $D$ lies in the kernel of the
linearisation is the order-one identity $P_kD_k+D_kP_k=D_k$, immediate from $w\in e$, $v\notin e$
and the corresponding statement for $f$, together with $D_e+D_f=0$. Finally the cross directions
are tangent, by the explicit rotations above, so the cone meets the kernel in more than the
origin while omitting $D$; and it is not closed under addition, since $Q$ is quadratic and does not
vanish on sums.
\end{proof}

The hypothesis asks only that some vertex lie in two hyperedges with a further vertex outside
both, which holds in every commuting tight family with $a\ge2$ and $n>|e\cup f|$; so the
singularity is not a feature of the two families measured. For the cross configuration the same
computation gives $\sigma_e(v)=-1$ and $\sigma_f(v)=+1$, which cancel, and that difference of sign
is the whole difference between the two kinds. Both the local computation and the full linear
system for $X$ are checked in \texttt{code/singular.py}, the latter by comparing the rank of the
coefficient matrix with that of the augmented matrix: $133$ against $134$ at Fano and $360$
against $361$ at $\mathrm{AG}(2,3)$ for the same-group direction, with equality in both for the
cross one. The scalar core is machine-checked as \texttt{TangentObstruction.no\_second\_order}.

That reduction is specific to projections, and the specificity is the point rather than a
caveat. Relaxing idempotency destroys it immediately, since $A_k=(b/p)I_p$ is positive semidefinite of
the right trace with $\sum_kA_k=aI$ and exceeds the band already at $p=4$, as recorded in
\S\ref{sec:leaves}. Its rank is $p$ rather than $b$, so what it shows is that trace and tightness
do not suffice; the rank-$n$ interpolation of \S\ref{sec:rank} is the statement that does keep the
rank fixed. That is not evidence against Xu's conjecture, which is about projections; it is the
obstruction recorded in \S\ref{sec:leaves} for $A_k=(b/p)I_p$, appearing again. It does mean any
argument for the monotonicity must use idempotency and not only the rank, the trace and
$\sum_kA_k=aI$, and it is why the search above enforces $\|A_k^2-A_k\|$ rather than assuming it:
without that check the retraction leaves the projection class silently, and the numbers it then
produces are about the wrong class.

The identification in Proposition~\ref{prop:bridge} is elementary and we make no claim of novelty
for it. For diagonal matrices it is stated without proof by Marcus, Spielman and Srivastava
\cite[\S7]{MSSII} in the square case $q=n$, with arbitrary diagonal weights; the statement here
covers $q\neq n$, specialises the weights to $0/1$ projections, supplies a proof, and carries the
substitution $x=t^2$ with its degree bookkeeping.
It should not be confused with the hypergraph Heilmann--Lieb theorem of Wan, Wang and Fan
\cite{Zhang}, which concerns a different polynomial --- the matching polynomial of the adjacency
hypermatrix, whose extreme root scales as $\Delta^{1/k}$ --- and neither implies the other. Both
routes of Proposition~\ref{prop:bridge} are computed in \texttt{code/xu\_sharp.py}, the mixed
characteristic polynomial from the definition
$\prod_e(1-\partial_{z_e})\det(xI+\sum_ez_eP_e)$ and the matching counts by subset dynamic
programming, sharing no code and agreeing coefficient by coefficient in exact arithmetic on four
hypergraphs at $b=2,3,4$. The skeleton of the argument is machine-checked: the root
identification of Proposition~\ref{prop:bridge} as \texttt{XuSharp.bridge\_roots}, the passage
from a ceiling and a converging ladder to the supremum as \texttt{XuSharp.sharp\_of\_ladder},
and the consequence that no smaller constant can hold as
\texttt{XuSharp.no\_smaller\_constant}. Godsil's bound and subgraph monotonicity enter there as
hypotheses, being classical.

Four further files carry the material of this subsection and of \S\ref{sec:mixeddisc}, all with no
\texttt{sorry} and on the three standard axioms. \texttt{RamaLean/SpectralNoGo.lean} is the
obstruction of \S\ref{sec:obstructions} in two halves that do not lean on each other:
\texttt{const\_of\_factors\_through\_traces}, that any functional reading only the power traces of
the blocks agrees on any two families of idempotents with equal block traces, which follows from
$A^{j+1}=A$ and nothing else; and \texttt{spectral\_bound\_is\_the\_target}, that a bound both valid
on the class and strong enough to give the target coincides with the target, with
\texttt{target\_gives\_spectral\_bound} for the converse. \texttt{RamaLean/ArcTangent.lean} carries
the identity $\sum_j\sigma_k(j)(D_k^2)_{jj}=0$ and the passage from a curve to Mathlib's
Mathlib's own tangent cone. \texttt{RamaLean/OrderFourSolvable.lean} is the step joining
\texttt{OrderFour.order\_four\_sum\_zero} to \texttt{CokernelRank.cokernel\_const}.
\texttt{RamaLean/MixedDiscriminant.lean} carries the symmetry, diagonal value and additivity of the
mixed discriminant. \texttt{RamaLean/ConeFix.lean} carries the entrywise solvability criterion behind
Proposition~\ref{prop:cone} and the identification of the obstructed pairs as the twin pairs, and
\texttt{RamaLean/ConeWitness.lean} the matrix level of the counterexample that makes the twin
conditions necessary, every step by kernel computation.
\texttt{RamaLean/OrderThreeLaw.lean} carries Lemma~\ref{lem:o3trace}.
\texttt{RamaLean/OrderThree.lean} is retained for its orbit lemmas; its
\texttt{diag\_mul\_add\_mul\_of\_cross} and \texttt{order\_three\_unobstructed} are true statements
whose hypothesis no second-order correction satisfies, and nothing here cites them. \texttt{RamaLean/SpanRank.lean} carries the two scalar criteria behind the span
and rank formulas, that a linear functional vanishes on $\{\sum t=0\}$ exactly when its coefficients
agree and that a diagonal quadratic form vanishes there exactly when every pair of coefficients
cancels, with the pigeonhole that makes $|K(i,j)|\ge3$ always an edge.
\texttt{RamaLean/TwoRegular.lean} carries every step of Proposition~\ref{prop:tworeg} except the
implicit function theorem itself: the divisibility, as the expansions of the quadratic and cubic
pieces at the jet $tD+t^2w$, where \texttt{quadratic\_expand} is the one place $Q(D)=0$ is consumed
and is what removes the $t^2$ term; that an affine map with surjective linear part has a zero, which
is the only place the rank hypothesis is used; and that a family of points of the set of the form
$x_0+tD+t^2w$ with $t\,w\to0$ puts $D$ in the tangent cone, which is where the second-order shape
does its work. What is cited and not formalised is stated in each file: the implicit function theorem
step of Proposition~\ref{prop:tworeg}, and Aleksandrov's, Bapat's and Gurvits' theorems.


\section{What a proof must contain}

\section{What a proof must contain}\label{sec:must}

Each of the methods that does not reach the band fails for a reason that can be
named exactly, and each such reason is a property that a proof must have. We
collect them here as a specification rather than as a list of failures: a proof
must not induct on vertex deletion, must not be a side-constant specialization of
the multivariate matching polynomial, must control deficient vertices of the path
tree, and must not pass through Floquet theory. Every one of these is sharp, in
the sense that the method in question is exactly optimal for the class of objects
it can see; that is the content of Section~\ref{sec:rank}.

The lower bound, unlike the upper bound, is not monotone
under vertex deletion. The Heilmann--Lieb argument for $|x|\le2\sqrt{\Delta-1}$
inducts on vertex deletion, which is available because $\Delta$ only decreases.
The quantity $|\sqrt{a-1}-\sqrt{b-1}|$ has no such monotonicity: deleting one
vertex of the larger side of $K_{3,4}$ produces $K_{3,3}$, which is regular, and
the gap collapses to zero, for $a=b$ one has $c=2\sqrt m$, so the interval
$[c-2\sqrt m,\,c+2\sqrt m]$ has its lower end at $0$ and the statement is empty. So the naive induction by vertex deletion, which is
what yields the upper bound, cannot prove the lower one.

That obstruction can be strengthened, and in a way that rules out the whole
family of such arguments rather than one instance of it. Every $b$-regular
bipartite graph is an induced subgraph of some $(a,b)$-biregular graph, by an
explicit padding construction; and over $b$-regular bipartite graphs the least
root of $\nu_G$ tends to $0$, in closed form $2\sin(\pi/2N)$ for $b=2$. So on the
class obtained by closing the biregular graphs under vertex deletion, the
infimum of the least root is $0$, and the inequality to be proved is simply
false there. Any argument that deletes vertices and applies an inductive
hypothesis, cavity recursions, stability transfer, the Heilmann--Lieb
induction, is an induction over that closure, and therefore cannot work,
not because biregularity is lost but because the statement itself fails. An induction that
carries something other than the target statement is not excluded, and \S\ref{sec:ratio} gives
one: its hypothesis $q-r\ge1$ is precisely a hypothesis the padding family above violates. For
$\mathrm{AG}(2,3)$, where $(s-t)^2=0.10102$, the minimum over all choices of $d$
deleted vertices on the larger side runs $0.3593$, $0.2424$, $0.1383$, $0.0549$
for $d=0,1,2,3$: it drops below the target at depth three.

A second route, substituting side-constant values into the multivariate matching
polynomial of Heilmann and Lieb, fails for a reason that is visible in one line.
Writing $M_G(z;x)=\sum_Mz^{|M|}\prod_{v\notin M}x_v$ and setting $x_v=\alpha$ on
the $p$-side and $x_v=\beta$ on the $q$-side gives
$\sum_km_kz^k\alpha^{p-k}\beta^{q-k}$, in which the exponent pair contributes
$c^{p-q}$ under $(\alpha,\beta)\mapsto(c\alpha,\beta/c)$ independently of $k$; so
the specialization depends on $\alpha$ and $\beta$ only through the product
$\alpha\beta$. Since the image of a product of two upper half-planes is
$\mathbb C\setminus[0,\infty)$, stability under this substitution says exactly
that the roots are real and nonnegative, and nothing more. In particular it says
the same thing when $a=b$, where the conclusion to be proved is empty. No choice
of $\alpha$ and $\beta$ can see $a-b$.

Two obstructions remain in the direct approach, both caused by cycles in $G$. First, a leaf of the path
tree that is an $A$-vertex carries $R_A=y>0$, of the wrong sign. This one is
disposed of: a maximal self-avoiding walk ending at $w\in A$ contains all $a$
neighbours of $w$, hence at least $a$ vertices of $B$ and therefore at least $a$
vertices of $A$, so
\[
p<a\ \Longrightarrow\ \text{the path tree has no }A\text{-leaf},
\]
which covers every $K_{p,q}$ with $p<q$. Second, and this is what blocks the
argument, an internal $A$-vertex of the path tree may have $k<P$ children, the
walk having already consumed part of its neighbourhood; the bound
$R_A\le y-P\alpha/(\alpha+Q)$ then fails. With $k=1$ closure would require
$\alpha(1-\alpha-Q)/(\alpha+Q)\ge y$, impossible for $b\ge3$ since $Q\ge2$.
Deficiency occurs already in $K_{3,4}$, where the walk $a_1b_1a_2b_2a_3$ leaves
$a_3$ with two children rather than three. What the argument as it stands
establishes is that a \emph{full} truncated biregular tree carries no eigenvalue
in the gap, consistent with the deficient subtree $P_8$ carrying one; that is a
statement about the tree and not about $G$. Controlling deficient $A$-vertices
is the open problem.

The case $b\ge3$ of the biregular family is the remaining open one, since $b=2$
is Theorem~\ref{thm:subdiv}; note that every graph named below has $a\ne b$, so
the statement being tested is not the empty one. It has been tested at $d=1$ on
$K_{3,4},K_{3,5},K_{3,6},K_{3,7},K_{3,8},K_{4,5},K_{4,6},K_{4,7},K_{5,6}$ and a
$(3,2)$ design incidence graph, with smallest-root ratios between $1.51$ and
$3.33$ and no root in a gap.

The natural attack is Godsil's recursion $R_v=x-\sum_{u}1/R_u$, whose fixed
points on the $(a,b)$-biregular tree are real precisely off the spectrum. Asking
that the region $R_b\in[x,P/x)$, $P=a-1$, $Q=b-1$, be closed under the
recursion reduces to $u^2-(2P+Q)u+P^2>0$ at $u=x^2$, i.e.\ to
$u<\bigl(2P+Q-\sqrt{Q^2+4PQ}\bigr)/2$. This is satisfied throughout the gap
with room to spare: the ratio of that threshold to $(\sqrt P-\sqrt Q)^2$ ranges
from $2.7$ to $37.7$ over the cases above. So the analytic part of an induction
is not the difficulty.

What blocks it is the boundary. A leaf of the path tree carries $R=x>0$, while
the argument wants $R<0$ at $a$-vertices, and mixed signs among the children of
a $b$-vertex leave the sum uncontrolled. One might hope to exploit that the
counterexample to the naive subtree argument, $P_8$ inside the $(3,2)$-biregular
tree, is deficient at \emph{internal} vertices, whereas a path tree is full at
internal vertices; but that is false for graphs of small girth, where a
self-avoiding walk can be blocked early, in $K_4$ the walk $0,1,2$ admits a
single extension. We therefore leave $b\ge3$ open.

A proof of Conjecture~\ref{conj:gap} cannot go through Floquet theory. For
$b_1=b$ the deck group of the universal cover is the free group $F_b$, and
Floquet theory is exactly the statement that for an \emph{abelian} deck group
the fibres are indexed by the characters of that group; the two coincide only
at $b=1$, where $F_1=\mathbb{Z}$. For $b\ge2$ the correct index set is the
unitary dual of $F_b$, and the relevant object is the adjacency operator in the
reduced $C^*$-algebra of $F_b$. This is not merely a technical obstruction: the
computations of \cite{PartI} show that the per-component weight in the
exponential formula, a single graph for $b=1$, becomes a sum over the many
non-isomorphic connected $\ell$-covers for $b\ge2$.

That framework is not empty. As $d\to\infty$, independent uniform random
permutations are strongly asymptotically free \cite{BC}, and consequently the
spectrum of a random $d$-lift of $G$ converges to $\operatorname{spec}(T)$.
Conjecture~\ref{conj:gap} is a finite-$d$ statement about the expected
characteristic polynomial rather than an asymptotic statement about typical
lifts, and its content is precisely that the containment holds uniformly in
$d$, with no exceptional set. What the $b_1=1$ case shows is that the interval
in \cite{HPS} is not the natural receptacle for the roots, and that the correct
one is spectral rather than metric.



\section{The formalization, in full}

\appendix

\section{The formalization}

The core algebraic derivation, the exponential formula, the closed form of
Theorem~\ref{thm:main} for all $n,r$ (equivalently, everything from the cycle
weights \eqref{eq:avg} to the closed form), its factorization, and the parity
Corollary~\ref{cor:parity}, is machine-checked in Lean~4 over Mathlib (v4.30).
The main entry points:

\begin{itemize}
\item \texttt{Paper2.expFormula}, the exponential formula
  \eqref{eq:newton} in integer form
  $r\,A_r=\sum_k r^{(k)} f(k)A_{r-k}$ over any commutative ring, proved via
  Mathlib's Cauchy cycle-type count
  (\texttt{Equiv.Perm.\allowbreak card\_of\_\allowbreak cycleType\_mul\_eq});
\item \texttt{Paper2.thm1\_Sr}, Theorem~\ref{thm:main} in the form:
  the expected cycle weight \eqref{eq:avg} equals the closed form, for all
  $n,r$, over any $\Q$-algebra domain;
\item \texttt{Paper2.cG\_factored}, \texttt{Paper2.quotient\_parity}, the factorization and Corollary~\ref{cor:parity}.
\end{itemize}

The three remaining corollaries are now formalized in full.
Corollary~\ref{cor:gf} is proved as a literal identity of formal power series,
\texttt{Paper2GenFun.genU\_mul} giving
$(1-2YX+X^2)\sum_dU_d(Y)X^d=1$ and \texttt{genPhi\_mul} the companion form with numerator
$(1-X)^2$, by coefficient comparison against the Chebyshev recurrence.
Corollary~\ref{cor:roots} is proved as the statement about roots,
\texttt{Paper2Roots.Phi\_ne\_zero\_of\_two\_lt\_abs}: if $|x|>2$ then
$\Phi_{n,r}(x)\neq0$, via $|T_n(x/2)|>1$ (\texttt{one\_lt\_abs\_cT}) and the vanishing of
neither factor of the factorization.
Corollary~\ref{cor:fib} is proved for all $n$ and $r$,
\texttt{Paper2FibLucas.cor2\_phi3\_general}, replacing the bounded
\texttt{native\_decide}; the route is the identity $F_{j+2m}+F_j=L_mF_{j+m}$ for even $m$
(\texttt{fib\_shift}), whose base cases are $F_{2m}=L_mF_m$ and, using Cassini,
$F_{2m+1}+1=L_mF_{m+1}$, and which is stated without truncated subtraction by writing
$L_m=2F_{m+1}-F_m$.

Trust base for the general theorems: \texttt{propext}, \texttt{Classical.choice},
\texttt{Quot.sound} only (checked with \texttt{\#print axioms}); the bounded
\texttt{native\_decide} checks additionally invoke \texttt{Lean.ofReduceBool}.
No \texttt{sorry}. The single remaining unformalized ingredient of the
\emph{graph-theoretic} statement is the sentence before \eqref{eq:avg}: a lift
of a cycle is a disjoint union of cycles, elementary covering-space
combinatorics for which Mathlib currently lacks the vocabulary.

For the biregular half, \texttt{RamaLean/BipartiteMatchingPoly.lean} supplies a
matching polynomial for bipartite graphs, Mathlib having none, and formalizes:
the deletion recursion $Q_G=Q_{G-v}-z\sum_{u\sim v}Q_{G-v-u}$
(\texttt{bipQ\_delete\_left}); the general Taylor shift
$[X^j]f(X+c)=\sum_i\binom{i+j}{j}f_{i+j}c^i$
(\texttt{taylor\_coeff\_of\_natDegree\_le}), from which the coefficient formulas at
$p-1$, $p-2$ and $p-4$ follow; the count $2m_2+|{\rm confL}|+|{\rm confR}|=|E|(|E|-1)$
with the two conflict classes as degree sums; and the combinatorial inputs to
Theorems~\ref{thm:universal} and~\ref{thm:c4}, that three pairwise-meeting edges
of a bipartite graph share a vertex (\texttt{three\_meeting\_share}), that the
conflict degree of an edge is $(d_u-1)+(d_v-1)$ (\texttt{conflict\_degree}), the
trichotomy for four-cycles of the conflict graph
(\texttt{four\_cycle\_trichotomy}), and the identity the whole expansion runs on,
that the matching indicator is the alternating sum over subsets of the conflicting
pairs (\texttt{alt\_sum\_conflictPairs}). Several of these, including the
trichotomy, depend on no axioms at all. \texttt{RamaLean/Paper2Unicyclic.lean}
further formalizes that the lower band edge is $(\sqrt P-\sqrt Q)^2$ and so
vanishes exactly when $P=Q$ (\texttt{gap\_eq\_zero\_iff}), which is the triviality
recorded above.

The inclusion--exclusion itself is formalized as a chain. The matching indicator is
the alternating sum over subsets of the conflicting pairs of a set of edges
(\texttt{alt\_sum\_conflictPairs}); substituting it turns $m_k$ into a double sum
(\texttt{mCount\_eq\_sum\_alt}); because conflict is intrinsic to a pair of edges,
the inner index set depends on the outer one only through a containment
(\texttt{powerset\_conflictPairs\_eq}), so the two summations exchange
(\texttt{mCount\_inclusion\_exclusion}); the resulting weight is a binomial
coefficient, by a bijection between the $k$-subsets containing a fixed $V$ and the
$(k-|V|)$-subsets of the complement
(\texttt{card\_powersetCard\_filter\_superset}), giving
\[
m_k=\sum_{S}(-1)^{|S|}\binom{|E|-|V(S)|}{k-|V(S)|}
\]
over sets $S$ of conflicting pairs (\texttt{mCount\_closed\_form}); and since the
weight sees $S$ only through $|S|$ and $|V(S)|$, the sum fibers over that pair
(\texttt{mCount\_grouped}, \texttt{mCount\_grouped\_range}), which is the grouping
by subgraph type. The chain is exercised end to end at $k=1$
(\texttt{mCount\_one\_via\_inclusion\_exclusion}): every nonempty set of
conflicting pairs has at least two endpoints, so only the empty set survives the
binomial and contributes $\binom{|E|}{1}$, reproving $m_1=|E|$ through the
inclusion--exclusion rather than by the direct bijection, so that two independent
routes cross-check each link.

What is \emph{not} formalized is the evaluation of the ten resulting counts as
explicit functions of $(p,a,b)$ and $C_4(G)$, the step at which
\texttt{three\_meeting\_share} and \texttt{four\_cycle\_trichotomy} are applied, together with the numerical verifications. One obstacle there is worth recording:
the formalized conflict set consists of \emph{ordered} pairs, so each edge of the
conflict graph occurs twice. This leaves the inclusion--exclusion untouched, since
the alternating sum vanishes over any nonempty index set whether or not it double
counts, but it means the grouping enumerates more configurations than the ten
unordered types above, and aligning the two would require rebuilding the conflict
set on unordered pairs.

The operator side of Sections~\cite{PartI}--\ref{sec:rank} is formalized in four
further files, all on the standard three axioms and with no \texttt{sorry}.
\texttt{RamaLean/SignReflection.lean} carries
Proposition~\ref{prop:reflect} to the point where only the rank-one expansion of
$\mu$ is left by hand: \texttt{sign\_avg\_reflect} proves that
$\sum_{\varepsilon\in\{\pm1\}^q}\det(tI-\sum_k\varepsilon_kS_k)$ is even or odd in
$t$ according to the parity of the dimension, with \emph{no} hypothesis on the
$S_k$, by the involution $\varepsilon\mapsto-\varepsilon$; \texttt{sum\_prod\_sgn}
is character orthogonality on the Boolean cube, proved by the same involution
applied to one coordinate; \texttt{sum\_multiaffine} deduces that averaging a
multi-affine function over signs returns $2^q$ times its constant term, which is
exactly what $\prod_k(1-\partial_{z_k})$ extracts; and
\texttt{mixedChar\_reflect} and \texttt{band\_edges\_swap} supply the shift and the
identity $2a-(\sqrt{a-1}+1)^2=(\sqrt{a-1}-1)^2$ that exchanges the band edges.
\texttt{RamaLean/ProductBound.lean} proves Proposition~\ref{prop:abbound}
(\texttt{no\_root\_above}: a convex combination of products of positive numbers is
positive) together with the exact overshoot $ab-(s+t)^2=(\sqrt m-1)^2$
(\texttt{edge\_gap}) and its vanishing exactly at $m=1$
(\texttt{edge\_gap\_eq\_zero\_iff}), which is the case $a=b=2$ where the band is
proved outright; it also proves Proposition~\ref{prop:ineq2}, that at $b=2$ the two
band edges are the single condition $\pi_{\max}\le4(a-1)/a^2$
(\texttt{band\_iff\_pi}), and that the strictly weaker $\pi_{\max}\le4/a$ already
yields $a+2\sqrt a$, which beats the Marchenko--Pastur edge for every $a$ and the
product bound $ab$ for $a>4$ (\texttt{weak\_bound}, \texttt{weak\_beats\_mp},
\texttt{weak\_beats\_ab}). \texttt{RamaLean/GramDet.lean} discharges the exterior-algebra inputs by writing them in
Gram coordinates, where each becomes a determinant statement. \texttt{det\_border} is the
bordered determinant formula
$\det\bigl[\begin{smallmatrix}a&r^{\mathsf T}\\ c&G\end{smallmatrix}\bigr]
=a\det G-r^{\mathsf T}\operatorname{adj}(G)c$, which is exactly \eqref{eq:wedgeip};
\texttt{det\_gram\_eq\_zero\_of\_dep} is $f_k\wedge\omega_k'=0$, the vanishing of a Gram
determinant on a dependent family; and \texttt{hsimple\_of\_border\_zero} combines them
into $\|f_k\|^2\|\omega_k'\|^2=\|\iota_{f_k}\omega_k'\|^2$, which is precisely the
hypothesis \texttt{hsimple} used below, so that hypothesis is now proved rather than
assumed. \texttt{RamaLean/Tightness.lean} supplies the other input, \texttt{htight}: that
$\operatorname{Adj}(A)=aI$ forces $\sum_k\iota_{f_k}\omega_k'=0$. Working in coordinates,
with a rank-two block presented by its two spanning vectors, the polarization step becomes
three elementary facts, \texttt{iota\_antisymm}, that
$\langle a,\iota_b\omega\rangle=-\langle b,\iota_a\omega\rangle$, which is a two-term
\texttt{ring} identity; \texttt{iota\_proj\_of\_orthogonal}, that compressing the block to
$e^\perp$ does not change the pairing against $f_k$; and that each $\iota_{f_k}\omega_k'$
already lies in $e^\perp$. \texttt{tight\_sum\_contraction\_eq\_zero} assembles them, and
\texttt{crossTerm\_nonneg\_of\_tight} composes the result with
\texttt{CrossTerm.crossTerm\_nonneg}, so that Theorem~\ref{thm:C2} is machine-checked from
the tightness hypothesis alone. \texttt{RamaLean/Plucker.lean} closes the last link, the identification of
\texttt{CrossTerm}'s coordinate vector \texttt{w} with the Pl\"ucker coordinates of
$\omega_k'$. Representing $b_1\wedge b_2$ by its antisymmetric array on \emph{ordered}
pairs scaled by $1/\sqrt2$ makes \texttt{pl\_dot}, that the resulting inner product is the
Gram determinant, hold with no $i<j$ bookkeeping: the double count and the two factors of
$1/\sqrt2$ cancel. \texttt{iota\_eq\_comb} writes $f_k=\iota_e\omega_k$ in the compressed
plane with explicit coefficients $-\langle e,b_2\rangle$ and $\langle e,b_1\rangle$, and
\texttt{hsimple\_of\_comb} then makes \texttt{hsimple} a \texttt{ring} identity rather
than an assumption. \texttt{crossTerm\_nonneg\_plucker} assembles everything:
\textbf{for a tight family of rank-two blocks and any unit vector, $C_2\ge0$ is
machine-checked with no hypothesis remaining}.

The same holds at every level. \texttt{GramDet.border\_gram} identifies the bordered Gram
matrix of $[v\mid M]$ with the Gram matrix of the bordered family;
\texttt{border\_det\_eq\_zero\_of\_mem\_span} then gives $f\wedge\omega=0$ whenever $f$
lies in the column span, which it always does; and \texttt{hsimple\_of\_mem\_span} combines
them into $\|f\|^2\|\omega\|^2=\|\iota_f\omega\|^2$ for a wedge of \emph{any} number of
vectors. So the hypothesis \texttt{hsimple} of Theorem~\ref{thm:Cr} is discharged at every
$r$, not only at $r=2$, and what Theorem~\ref{thm:Cr} asserts, that each $S$-summand of
$C_r$ is a Schur square minus the square of a contraction, is machine-checked
unconditionally. Finally \texttt{gram\_det\_basis\_of\_injective} and
\texttt{gram\_det\_basis\_of\_not\_injective} prove that for a coordinate family the Gram
determinant of a set of blocks is $1$ on matchings and $0$ otherwise, that is
$F_A=\mu_G$; that is the first of the two inputs to Proposition~\ref{prop:km}, and the
second, the Kesten--McKay limit, is the only analytic ingredient anywhere in this section
and is cited rather than proved.
\texttt{RamaLean/CrossTerm.lean} carries Theorems~\ref{thm:C2} and~\ref{thm:Cr}.
\texttt{crossTerm\_eq\_sq\_sub} is the general level: with no hypothesis beyond
\texttt{hsimple} it proves
$\sum_{k\ne l}[\langle f_k,f_l\rangle\langle w_k,w_l\rangle-\langle u_k,u_l\rangle]
=\|\sum_kf_k\otimes w_k\|^2-\|\sum_ku_k\|^2$, which is the $S$-summand of
Theorem~\ref{thm:Cr} once $w_k=\omega_k'\wedge\omega_S'$ and
$u_k=\iota_{f_k}(\omega_k'\wedge\omega_S')$. Its two geometric
inputs are hypotheses, \texttt{hsimple}, that $f_k\wedge\omega_k'=0$ in its metric form
$\|f_k\|^2\|\omega_k'\|^2=\|u_k\|^2$, and \texttt{htight}, that $\sum_ku_k=0$, and
\texttt{crossTerm\_eq\_sq} specializes it to $r=2$, where tightness kills the second
square, and proves that they force $C_2$ to equal
$\sum_{i,j}\bigl(\sum_k(f_k)_i(\omega_k')_j\bigr)^2$, with \texttt{crossTerm\_nonneg} and
\texttt{crossTerm\_eq\_zero\_iff} the consequences. The supporting
\texttt{sum\_gram\_prod\_eq\_sum\_sq} is the Schur product theorem in the case needed,
that $\sum_{k,l}\langle f_k,f_l\rangle\langle w_k,w_l\rangle$ is a sum of squares. No
exterior algebra is involved: the bivectors enter only through their coordinates.
\texttt{RamaLean/VertexSplit.lean} carries the Gram form of
Theorem~\ref{thm:vertexrec}. Its core, \texttt{det\_add\_vecMulVec}, is the adjugate
form \eqref{eq:mdl} of the matrix determinant lemma; Mathlib has the lemma only as
$\det(A+uv^{\mathsf T})=\det A\cdot\det(1+v^{\mathsf T}A^{-1}u)$ and records the
adjugate form as an explicit \texttt{TODO}, so this supplies it, under the hypothesis
that $\det A$ is a unit, which is what the application needs, the Gram matrix of an
independent family being positive definite. \texttt{det\_le\_det\_add\_vecMulVec} is
the inequality the recursion consumes, $M_r(A)\ge M_r(A^{(e)})$, deduced from
\texttt{AdjugatePSD.adjugate\_posDef}. The singular case is scoped out as in
\texttt{AdjugatePSD}: both sides of \eqref{eq:mdl} are polynomial in the entries, so it
extends by continuity from $A+\varepsilon I$, and that limit is not formalized.
\texttt{RamaLean/CavityThreshold.lean} carries
Proposition~\ref{prop:cross}: \texttt{cavity\_step} is the induction step, that
child ratios at least $x/2$, nonnegative weights of total at most $a$, a nonpositive
remainder and $4a\le x^2$ force the parent ratio to be at least $x/2$ again;
\texttt{threshold\_iff} is its sharpness, that $x-2a/x\ge x/2$ holds \emph{if and
only if} $x\ge2\sqrt a$; \texttt{double\_root\_at\_threshold} is the structural
form of the same statement, the factorization
$\rho^2-2\sqrt a\,\rho+a=(\rho-\sqrt a)^2$ that collapses the interval of admissible
thresholds to a point; and \texttt{no\_slack\_propagates} is
Proposition~\ref{prop:rigid}; \texttt{cavity\_step\_deficient} is
Proposition~\ref{prop:defstep}, the strengthened step; \texttt{compression\_sum},
\texttt{trace\_le\_of\_rayleigh} and \texttt{remainder\_upper} are
Proposition~\ref{prop:twolevel}; and \texttt{ratio\_form},
\texttt{one\_add\_le\_inv\_one\_sub}, \texttt{inv\_one\_sub\_sub\_one\_add},
\texttt{km\_slack\_calc} and \texttt{km\_satisfies\_target} are
\eqref{eq:ratiotarget} and Proposition~\ref{prop:kmtarget}. Proposition~\ref{prop:deficient} is
\texttt{Tightness.adj\_compressed}, with the per-block splitting
\texttt{Tightness.iota\_split}. \texttt{RamaLean/KestenMcKay.lean} carries
Proposition~\ref{prop:km}: \texttt{radical\_at\_edge} that the radical is exactly $2$ at
the threshold, \texttt{G\_at\_edge} the closed form there, \texttt{inv\_G\_div\_edge} the
constant $1+1/(1+\sqrt a)$, and \texttt{excess\_eq}, \texttt{excess\_pos},
\texttt{excess\_lt} that the excess is exactly $1/(1+\sqrt a)$, positive for every $a>1$
and below any $\varepsilon$ once $\sqrt a>1/\varepsilon$. The two inputs it does not
formalize are the ones that are not arithmetic: that $F_A$ is the matching polynomial,
which is Theorem~\ref{thm:matchingform}. The Kesten--McKay step is no longer cited as a
limit: \texttt{RamaLean/MomentTransfer.lean} carries Proposition~\ref{prop:transfer} in
full, namely the exact geometric expansion with remainder
(\texttt{geom\_remainder}), the remainder bound (\texttt{remainder\_le}), the moment
transfer (\texttt{stieltjes\_close}), its form for reciprocals (\texttt{inv\_close}) and the
positivity of the gap between the support and the evaluation point
(\texttt{edge\_gap\_pos}). What remains cited is the locality of root moments, which is
\cite{CsF} after \cite{AH}, and Godsil's tree-like walk formulation \cite{Go}. The vertex recursion of Theorem~\ref{thm:vertexrec} and the
sign of $X_e$ enter as hypotheses; what is machine-checked is that they close the
induction, and that they close it exactly at $2\sqrt a$.
\texttt{RamaLean/TracePSD.lean} and
\texttt{RamaLean/AdjugatePSD.lean} carry the two positivity facts the extremal
computation and the kernel recursion run on. In both, the standard input is a
visible hypothesis rather than a hidden one: \texttt{trace\_sq\_le\_sq\_trace}
takes the spectral decomposition of a positive semidefinite matrix as an explicit
assumption, and \texttt{adjugate\_posDef} is proved in the definite case, the
semidefinite limit $A+\varepsilon I$ being left by hand.

\texttt{RamaLean/LaguerreBand.lean} carries Theorem~\ref{thm:kpq} in full on the
inequality side: \texttt{sqrt\_amgm} and \texttt{sqrt\_step} are the arithmetic--geometric
mean step \eqref{eq:amgm}, \texttt{rowRadius\_le} and \texttt{lastRow\_le} are the two
Gershgorin row bounds, \texttt{band\_endpoints} identifies the disc with the biregular
tree band, and \texttt{laguerre\_band} assembles them through Mathlib's
\texttt{eigenvalue\_mem\_ball}. What is not formalized is the classical passage from
$\mu_{K_{p,q}}$ to the Laguerre polynomial and thence to its Jacobi matrix; that
identification is a coefficient comparison and is done by hand above.

The three developments of \cite{PartI} are formalized to the extent stated
there. \texttt{RamaLean/AbelianCover.lean} proves that the integrand is real
(\texttt{det\_sub\_smul\_one\_im}), that a continuous real function with vanishing average
on a connected space has a zero (\texttt{exists\_zero\_of\_integral\_zero}), and the
resulting localization and reduction; the Godsil--Gutman expansion is assumed as the
hypothesis \texttt{haverage}. \texttt{RamaLean/TreeSubstitution.lean} proves the
substitution identity as a homogenization (\texttt{homogenize}), fixes its orientation on
the star (\texttt{star\_ratio}), and transports Heilmann--Lieb through a subdivision
(\texttt{subdivision\_root}, \texttt{subdivision\_abs}); the covering combinatorics is by
hand. \texttt{RamaLean/FeedbackVertex.lean} proves the combinatorial core
(\texttt{no\_bypass}, via Mathlib's characterization of acyclicity by path uniqueness), the
resolvent diagonal, the algebraic collapse of the trace formula, and the trace step; the
Pimsner--Voiculescu input appears as an explicit hypothesis of
\texttt{feedback\_one\_of}, since it is not in Mathlib. Every theorem in all four files
depends only on the three standard axioms, and none uses \texttt{native\_decide}.

\begin{remark}
Everything here is stated for the uniform permutation model; by
\cite{HPS} the quotient coincides with the $(r-1)$-matching polynomial,
so Theorem~\ref{thm:main} can alternatively be deduced from \cite{CGHRS}, with the caveat
about that preprint's status recorded in \cite{PartI}'s preamble.
We believe the generating function of Corollary~\ref{cor:gf}, the
evaluation of Corollary~\ref{cor:fib}, and the formalization are new.
\end{remark}




\section{Saturation, the path tree, and the free-convolution reading}

\subsection{The bound is per-vertex, and the saturation is only of constants}\label{sec:pervertex}

The limit at $(3,6,5)$ is real but narrower than it looks. What is saturated there is the
\emph{uniform} certificate: both bounds are attained with ratio $1.0000$, so no choice of
constants does better (\texttt{saturation\_blocks\_improvement}). The refinement below is not a
choice of constants, and it escapes for that reason.

The requirement \eqref{eq:closure} caps every child by the same $\kappa$, evaluated at the worst
child count $d-1$. But Proposition~\ref{prop:ratiosteps}(i) already gives a left-type vertex with
$j$ children the sharper bound $\kappa(j)=\lambda+j/m$, and a right-type vertex with very few
children sits where the path tree is nearly exhausted, so its children are child-poor too and
carry a much smaller $\kappa$. Instantiating the bound at each child's own count rather than
globally, the requirement at a right-type vertex becomes
\[
   k\;>\;\lambda\,\kappa(j_{\max}),\qquad j_{\max}=\max_i j_i ,
\]
which is strictly weaker whenever $j_{\max}<d-1$ (\texttt{RatioRoute.refined\_weaker},
\texttt{right\_step\_refined}). Measured on the path trees, it fails nowhere, including at the
two $(3,6,5)$ points where the uniform requirement fails at eight vertices
(\texttt{code/twolevel.py}). The saturation result is untouched: it says the constants cannot be
improved, and this does not improve them.

Substituting Lemma~\ref{lem:blocking} for $j_{\max}$ and asking the requirement for every
admissible $k$ removes the path tree altogether. With $\Delta=q-r$ it reduces, by one
cancellation, to two inequalities:

\begin{theorem}\label{thm:structclose}
If $\lambda<m$ and $\lambda^2<\Delta$, then $\lambda\,\kappa(j)<k$ for every $k\ge\Delta$ and
every $j\le k-\Delta$.
\end{theorem}

\begin{proof}
$\Delta(m-\lambda)>\lambda(\lambda m-\Delta)$ is equivalent to $\Delta>\lambda^2$ after
cancelling $m$; since $m>\lambda$ and $k\ge\Delta$ the left side only grows with $k$
(\texttt{RatioRoute.structural\_closure}).
\end{proof}

Neither hypothesis mentions a path tree, a family or a graph. The second is the binding case
already noted, where the children are leaves of ratio exactly $\lambda$; the first is what the
per-vertex bound buys. Over $354$ parameter points $(d,q,r)$ that are \emph{realizable} --- a
$(d,q)$-biregular graph with $r$ vertices of degree $q$ needs $qr/d$ to be a positive integer ---
the pair holds at $326$, or $92.1\%$, failing only as $\lambda$ approaches the gap edge
(\texttt{code/universal\_close.py}). The elementary sharpening $k_u\le\min(d-1,k-\Delta)$, using
that a child lies in $L$ and has degree $d$, gains nothing at any of them
(\texttt{code/sharp\_close.py}): the obstruction is the two conditions themselves and not the
counting.

\subsection{The path tree inherits the gap}\label{sec:inherit}

That obstruction is the invariant's and not the method's, and the reason is worth recording
because it says the route is not exhausted. The recursion computes $F_v=\mu_{T_v}/\mu_{T_v-v}$,
so ``no $F$ vanishes'' is exactly $\mu_P(\lambda)\ne0$ --- strictly stronger than
$\mu_G(\lambda)\ne0$, since $\mu_P$ carries its own roots as well as those of $\mu_G$. Were
$\mu_P$ to vanish somewhere in the gap, no certificate of this kind could succeed there,
whatever invariant it used.

It does not. On a tree $\mu_P=\prod_vF_v$ and the $F_v$ are the pivots of the $LDL$ factorization
of $\lambda I-A_P$, so by Sylvester's law the number of negative $F_v$ counts the eigenvalues of
$P$ above $\lambda$; comparing the count at the two ends of the gap counts the eigenvalues
between, in time linear in $|P|$ and with no matrix. Over four $(d,q,r)$ families that are
\emph{not} complete bipartite, with path trees of $3457$, $17241$, $22402$ and $291429$ vertices,
the count is equal at both ends in every case: the path tree has no eigenvalue in the gap
(\texttt{code/pathtree\_inertia.py}). This is `HEURISTIC', four families, and it is stronger than
Godsil's divisibility, which places the roots of $\mu_G$ among those of $\mu_P$ and says nothing
about where the others go.

So the statement the route proves is true throughout the gap, including where the
sign-alternating certificate cannot hold. An argument therefore exists in principle at every
$\lambda$ in the gap. What does not exist is a better certificate of \emph{this} family, and that
is worth separating from the first point.

Two things settle it. Measuring $\min_v|F_v|$ across the gap suggests the sign-free alternative
$|F_v|\ge\lambda$: the minimum equals $\lambda$ exactly, attained at a leaf, at every fraction
tested in two of three families, falling to $0.83\lambda$ and $0.51\lambda$ only at $(3,6,5)$ near
the edge (\texttt{code/minF\_scan.py}). But it is the harder condition, not a weaker one. With
every child positive, $F_v=\lambda-\sum_u1/F_u$ falls below $\lambda$ automatically, and
$|F_v|\ge\lambda$ additionally forces the sum past $2\lambda$, which at a right-type vertex is
$k\gtrsim2\lambda B$ against the $k>\lambda B$ the sign condition needs --- a factor of two in the
child count.

And there is no cancellation to repair. The path tree is bipartite by depth, so the children of a
vertex all lie on one side and are of one type, and same-type siblings never disagree in sign: over
three families at six gap fractions each, including $0.99$, the number of vertices whose children
have mixed signs is \emph{zero} in every case, out of $865$, $4185$ and $3220$ parents
(\texttt{code/signmix.py}).

\subsection{The same measurement reformulates the conjecture}\label{sec:reform}

The inheritance of \S\ref{sec:inherit} says more than that the ratio route's target is true. It
turns the conjecture into a statement about a finite tree.

Godsil gives $\mu_G\mid\mu_P$, and $P$ is a tree, where the matching polynomial \emph{is} the
characteristic polynomial. So every root of $\mu_G$ is an eigenvalue of $P$, and

\begin{proposition}\label{prop:reform}
$\operatorname{spec}(P)\subseteq\operatorname{spec}(T)$ implies
$\operatorname{Zeros}(\mu_G)\subseteq\operatorname{spec}(T)$. Conversely a root outside
$\operatorname{spec}(T)$ is an eigenvalue of $P$ outside it, so the containment fails exactly
where the conjecture does.
\end{proposition}

\begin{proof}
Immediate from $\mu_G\mid\mu_P$ in both directions
(\texttt{PathTreeRoute.conj\_of\_pathtree}, \texttt{pathtree\_fails\_of\_violation}).
\end{proof}

The hypothesis is strictly stronger than the conclusion, since $\mu_P$ carries its own roots as
well, and that is the point rather than a defect: it is a statement about the spectrum of one
explicit finite tree with no matchings in it. It is also not vacuous in either direction. It fails
for Hall's graph, necessarily --- $\sqrt5$ is a root of $\mu_G$, hence an eigenvalue of $P$, and
the ratio system at $\sqrt5$ has decay $0.9636233789$, so it is outside $\operatorname{spec}(T)$
--- and a hypothesis provable in general would be worthless, the conjecture being false. And it
holds wherever it has been measured, over the four families of \S\ref{sec:inherit}.

The strictness of the hypothesis is decisive, and against it. At minimum degree three the
containment is \emph{false}. Over $86$ graphs with $\delta\ge3$ carrying a gap of
$\operatorname{spec}(T)$, all but one inherit it; the exception is a $12$-vertex graph with
degrees $\{3,4,10\}$ whose path tree, on $3944$ vertices, carries twelve eigenvalues inside the
gap $[2.060,2.160]$ while no root of $\mu_G$ lies there --- counted by Sylvester's law, with the
near-zero-pivot flag honoured, so the count is trustworthy
(\texttt{code/pathtree\_delta3.py}).

So $\operatorname{spec}(P)\subseteq\operatorname{spec}(T)$ fails exactly where
Conjecture~\ref{conj:gap} holds, and Proposition~\ref{prop:reform} cannot carry a proof of the
minimum-degree conjecture: no argument can pass through a hypothesis that is false on the graphs
in question. What survives is the biregular reading. There the containment is measured to hold
(\S\ref{sec:inherit}), and there the proposition does reduce the inner end to a statement about
one explicit finite tree. We record the negative result because it closes a route that looks
plausible from the biregular evidence alone, and because it was the only reformulation five
searches produced.

So the sign-alternating certificate is the best of its family, and its reach is exactly
$\lambda^2<\Delta$ --- the binding case, where the children at the minimum count are leaves of
ratio $\lambda$ and the parent is $\lambda-\Delta/\lambda$. \textbf{That boundary is sharp rather
than an artefact of the certificate's shape}, and extending the route past it needs an argument
that is not an interval invariant on the ratios at all, since the target it would prove is true
there and every certificate of this kind provably is not.

Finally, the sign separation is a consequence of bipartiteness: on wheels no natural class
determines the sign of a path-tree ratio, endpoint degree, depth parity and child count each
being mixed (\texttt{code/d3ratio.py}), so the certificate does not transfer to
Conjecture~\ref{conj:mindeg3} in general.


The biregular case is what survives, and \S\ref{sec:ratio} settles it on a nontrivial class by
the ratio recursion. This part develops a second route to the same statement, through counting,
the coefficient ladder and the operator form, together with the obstructions each of them meets.
It uses the framework of Part~I throughout; nothing in Part~I depends on it, except that
\S\ref{sec:ratio} and \cite{PartI} each cite one result proved here.

\section{The band is the support of a free convolution power}

There is a single reason behind all of this, and it is worth stating because it
accounts uniformly for the failure of every method tried. Let
$\chi=(1-\tfrac1b)\delta_0+\tfrac1b\delta_b$, and let $\nu$ be the centred slot
spectrum $\tfrac1b\delta_{b-1}+(1-\tfrac1b)\delta_{-1}$, so that
$\chi=\delta_1\boxplus\nu$.

\begin{proposition}\label{prop:freepower}
$\operatorname{supp}\bigl(\chi^{\boxplus a}\bigr)
=\bigl[(\sqrt{a-1}-\sqrt{b-1})^2,\ (\sqrt{a-1}+\sqrt{b-1})^2\bigr]$.
\end{proposition}

\begin{proof}
$\chi$ is the spectral distribution of a rank-$b$ projection scaled so that $a$ free copies sum to
the identity, and the free additive convolution power of such a measure is the spectral measure of
the $(a,b)$-biregular tree, whose spectral radius is $\sqrt{a-1}+\sqrt{b-1}$ and whose spectrum is
the stated interval; this is the same tree computation used in the proof of
Corollary~\ref{cor:xucoord}. It is classical and is used here only to name the band; no novelty is
claimed,
and nothing below depends on it beyond the identification of the two endpoints already established
in Proposition~\ref{prop:bridge} and Corollary~\ref{cor:xucoord}.
\end{proof}

So the band is the support of a free convolution \emph{power}. Every method that
fails computes instead a compound free Poisson of rate $a$ with jump $\nu$, whose
free cumulants are $a\,m_n(\nu)$ rather than $a\,\kappa_n(\nu)$, and
\[
m_n(\nu)=\kappa_n(\nu)\ \ (n\le3),\qquad
\kappa_4(\nu)-m_4(\nu)=-2(b-1)^2 .
\]
The two sequences agree through third order and part company at the fourth, by
exactly $-2a(b-1)^2$. Order four is where every structure in this problem turns
over, and it does so three independent times: the coefficients $E_0,\dots,E_3$ are
determined by the parameters while $E_4$ carries the $4$-cycle count; the moment and
free-cumulant sequences of the jump agree up to order three and split at order four;
and the parallel-class construction for the finite free convolution reproduces the
elementary symmetric functions $e_1,e_2,e_3$ of $\mu$ and first fails at $e_4$. That is why the barrier bound, the scalar free convolution
and the finite free convolution all overshoot by the same margin of about $1.1$,
and it is the analytic form of the third-moment obstruction above.

The finite free cumulants of $\mu$ itself follow the free power, not the Poisson:
\[
\kappa_1=a,\qquad \kappa_2=\frac{ap(b-1)}{p-1},\qquad
\kappa_3=\frac{ap^2(b-1)(b-2)}{(p-1)(p-2)},
\]
with $\kappa_4$ the first that depends on the family, matching the coefficient
ladder, where $E_0,\dots,E_3$ are universal and $E_4$ carries the $4$-cycle count.
Note $\kappa_3=0$ exactly when $b=2$: Proposition~\ref{prop:reflect} is visible in
the cumulants.

Two further facts delimit what can be done with this. Writing
$\psi_0(x)=x^{p-p/b}(x-b)^{p/b}$ and $\Psi=\psi_0^{\boxplus_pa}$, the polynomial
$\Psi$ satisfies the band at every $p$, but $\mu\ne\Psi$ in general and there is no
real-rooted $\rho$ with $\mu=\Psi\boxplus_p\rho$ unless the two agree: finite free
deconvolution is unique, so such a $\rho$ would have $\kappa_1=\kappa_2=0$, and
$\kappa_2$ is a positive multiple of the variance of the roots, forcing
$\rho=x^p$. The very agreement of the first three cumulants is what forbids the
factorization. Divisibility by a single copy, $\mu=\psi_0\boxplus_p\rho$ with $\rho$ real-rooted,
would by the Marcus--Spielman--Srivastava bound give a lower bound on
$\lambda_{\min}$, and it holds in small cases; but it is not a property of the
class. Taking $b=2$ and the incidence families of $3$-regular graphs, the
deconvolution $\rho$ is real-rooted for $S(K_4)$ at $p=4$ and for $K_{3,3}$ at
$p=6$, and fails for the cube $Q_3$ at $p=8$ and the Petersen graph at $p=10$,
where the largest imaginary part of a root of $\rho$ is $0.133$ and $0.290$
respectively. The cube is the informative case: it is $3$-regular and bipartite,
hence $1$-factorable, so the failure is not caused by the absence of a partition of
the blocks into parallel classes. Whatever separates the projection class from the
positive semidefinite relaxation, single-copy divisibility is not it: at
$(p,a,b)=(12,4,2)$ the scalar family, which violates the band, admits the
factorization while a projection family does not.

Two general facts about free convolution came out of the attempt, and they bound
what any argument of this shape can achieve. For finitely supported probability
measures,
\begin{equation}\label{eq:squeeze}
\min\operatorname{supp}\omega+\min\operatorname{supp}\tau
\ \le\ \min\operatorname{supp}(\omega\boxplus\tau)
\ \le\ \min\operatorname{supp}\tau+m_1(\omega),
\end{equation}
the lower bound from the operator model and the upper from
$\min\operatorname{supp}(\omega\boxplus\tau)=\sup_{w<0}[K_\tau(w)+R_\omega(w)]$
together with the monotonicity of $R_\omega$, which follows from Cauchy--Schwarz.
And if $\min\operatorname{supp}\tau\ge c$ then
$\tau(\{c\})\le\mu_2/\bigl(\mu_2+(m_1-c)^2\bigr)$, so by Bercovici--Voiculescu an
atom of mass exceeding $\omega(\{\min\operatorname{supp}\omega\})$ pins the left
edge of $\omega\boxplus\tau$ at $c$ exactly.

Applied with $\omega=\chi$, for which $m_1=1$, \eqref{eq:squeeze} says that
convolving with $\chi$ moves the left edge by at most $1$. Any certificate that
tries to reach the band edge by exhibiting $\mu$ as a free convolution with $\chi$
is therefore capped, and the first three cumulants do not constrain the other
factor enough: there are measures with exactly the forced values of
$\kappa_1,\kappa_2,\kappa_3$, real-rooted, supported in $[0,ab]$, whose convolution
with $\chi$ has left edge $0$. At $(a,b)=(5,4)$, $p=12$ one may take
$x^4(x-4)^3(x-6)^4(x-12)$, whose mass at $0$ exceeds $1/b$.

The same atom mechanism explains the hypothesis $a\ge b$: for $a<b$ the measure
$\chi^{\boxplus a}$ carries an atom at $0$ of mass $1-a/b$, so its left edge is $0$
and the band statement is empty.

\section{Rank monotonicity, and what it explains}\label{sec:rank}

What does separate them can be said exactly, in terms of the representation
\eqref{eq:dpp}. For a general positive semidefinite family with $\sum_kA_k=aI$ the
marginal of $s_k$ has probability generating function
$\prod_i\bigl(\tfrac{a-\alpha_i}{a}+\tfrac{\alpha_i}{a}z\bigr)$ over the eigenvalues
$\alpha_i$ of $A_k$, a Poisson binomial; it is exactly $\mathrm{Binomial}(b,1/a)$ if
and only if every $\alpha_i$ lies in $\{0,1\}$, that is if and only if $A_k$ is a
rank-$b$ projection. Among Poisson binomials of fixed mean $b/a$ with $n$ equal
parameters the variance is $(b/a)(1-b/(na))$, increasing in $n$; so the projections,
where $n=b$, are the \emph{variance-minimizing} extreme point of the whole family.
The scalar counterexample $A_k=(b/p)I$ sits at the opposite end, at the Poisson
limit, with variance exceeding the projection value by exactly $b/a^2$.

This has a consequence that organizes the whole picture. Fix $(a,b)$ and let the
common rank $n$ of the $A_k$ vary, with $A_k=(b/n)Q_k$ for rank-$n$ projections
$Q_k$ summing to $(an/b)I$, so that $\operatorname{tr}A_k=b$ and $\sum_kA_k=aI$
throughout. At $(a,b)=(4,2)$, $p=6$, $q=12$ the extreme roots move monotonically:
\begin{center}
\begin{tabular}{ccc}
\hline
$n$ & least root & greatest root\\
\hline
$2$ & $1.129$ & $6.871$\\
$3$ & $1.004$ & $7.706$\\
$4$ & $0.951$ & $8.090$\\
$5$ & $0.921$ & $8.317$\\
$6$ & $0.903$ & $8.467$\\
\hline
\end{tabular}
\end{center}
against the band $[0.536,7.464]$ at $n=b=2$ and the Marchenko--Pastur interval
$[0.343,11.657]$ approached as $n\to p$. So the conjectured band is the endpoint at
$n=b$ of a family of intervals increasing in $n$, whose other endpoint is
Marchenko--Pastur, and the Marchenko--Pastur value is attained. That is the correct
reading of the obstructions recorded above: each of them is sharp for the class it
sees, namely the one with $n$ unrestricted, and the conjecture is the assertion that
restricting to $n=b$ contracts the interval all the way to the tree band.

The condition $\operatorname{rank}A_k\le b$ is therefore the condition that the
occupancy of a block fluctuates as little as possible at its given mean, and the
conjecture asserts that minimal marginal fluctuation propagates to the extreme
confinement recorded above. That suggests the right statement is one about measures
rather than matrices: for a random vector $(s_1,\dots,s_q)$ of nonnegative integers
with $\sum_ks_k=p$ almost surely, each $s_k$ exactly $\mathrm{Binomial}(b,1/a)$, and
joint generating function homogeneous of degree $p$ and real stable, are the roots
of $\mathbb E\prod_k(y-as_k)$ confined to the band? Every projection family induces
such a law, and the scalar family does not, failing the marginal condition.

That class is rigid, which is what makes the reformulation faithful rather than a
weakening. The marginal condition is itself a rigidity statement: in the polarized
picture it holds if and only if the $b$ slots inside each block are exactly jointly
independent Bernoulli$(1/a)$. Consequently, any \emph{determinantal} law satisfying
the two combinatorial conditions has a rank-$p$ projection kernel whose diagonal
$b\times b$ blocks are $\tfrac1aI_b$, that is, it is the Naimark kernel of a
projection family; and any law given by independent balls in bins satisfying them
comes from an $(a,b)$-biregular incidence matrix, that is, from a commuting family.
So the two canonical constructions of stable laws produce nothing new. At $b=2$ more
is true: for $(p,q,a,b)=(4,6,3,2)$ and $(3,6,4,2)$ the affine hull of the stable laws
in the class coincides with the affine hull of the projection families, of dimensions
$54$ and $14$ against $78$ and $38$ for the two combinatorial conditions alone. The
measure formulation is therefore not a relaxation at $b=2$.

The two combinatorial conditions alone do not suffice, and the point at which they
fail is sharp: over all laws satisfying them the largest root is confined to the band
only for $q$ below an explicit threshold, and it reaches the trivial value $ab$
exactly at $q=a^b$, the smallest $q$ at which the binomial profile
$n_j=q\binom bj(a-1)^{b-j}a^{-b}$ is integral. Stability is what must supply the
rest, and it cannot be seen by any argument that symmetrizes: the divisibility
constraint that stability imposes, namely that $(w+a-1)^{b-1}$ divides
$\mathbb E\bigl[s_l\,w^{s_k}\bigr]$, becomes vacuous on exchangeable laws.
Negative association alone is likewise insufficient, by an explicit permutation
counterexample whose polynomial has a root at $ab$.

The size of the claim is also worth recording. Since $E_1$ and $E_2$ are universal,
the root distribution of $\mu$ has mean $a$ and variance $a(b-1)$ for every
admissible family. Measured against that, the lower band edge sits between $1.35$
and $1.92$ standard deviations below the mean over the range $3\le a\le100$,
tending to $2$ as $a$ grows. The conjecture is therefore a statement of extreme
concentration: the least of $p$ roots is confined to about two standard deviations,
where an unstructured sample of that size would reach out to order
$\sqrt{\log p}$.

\subsection{The extremal problem in dimension four}

Graphs are not always extremal in that class, however, and the smallest case can be
settled exactly. This is no longer needed for the band, at $p=4$, $b=2$ the
reflection gives $\mu(t+a)=t^4-2at^2+\beta$, and real-rootedness alone forces
$0\le\beta\le a^2$, hence $t^2\le2a\le4(a-1)$, but it identifies the extremal
families. Take $b=2$, $p=4$, so $q=2a$ and $\sum_kP_k=aI_4$. Writing
$g_{ij}=\operatorname{tr}(P_iP_j)$ and $h_{ij}=\operatorname{tr}\bigl((P_iP_j)^2\bigr)$,

\begin{proposition}\label{prop:E4}
For rank-$2$ orthogonal projections $P_1,\dots,P_q$ on $\mathbb R^4$ with
$\sum_kP_k=\tfrac q2I$,
\[
E_4=\frac{q^4}{16}-\frac{q^3}{4}+\frac q2+\frac14\sum_{i\ne j}\bigl(g_{ij}^2-h_{ij}\bigr),
\]
and every summand is nonnegative. Equality holds exactly when
$\operatorname{rank}(P_iP_j)\le1$ for all $i\ne j$.
\end{proposition}

Nonnegativity is immediate once the right object is named: for $i\ne j$ put
$X=P_jP_iP_j\succeq0$; then $\operatorname{tr}\bigl((P_iP_j)^m\bigr)=\operatorname{tr}(X^m)$
for every $m\ge1$, so $g_{ij}^2-h_{ij}=(\operatorname{tr}X)^2-\operatorname{tr}X^2
=2\sum_{r<s}\lambda_r\lambda_s\ge0$. The last step is formalized in
\texttt{RamaLean/TracePSD.lean}: \texttt{sq\_sum\_sub\_sum\_sq} gives the identity
$\bigl(\sum f\bigr)^2-\sum f^2=\sum_i\sum_{j\ne i}f_if_j$ over any finite index set
and \texttt{sum\_sq\_le\_sq\_sum} the inequality under nonnegativity, with the
spectral input carried as explicit hypotheses on the matrix form.

Identifying a $2$-plane $U\subset\mathbb R^4$ with its unit simple bivector
$\omega_U\in\Lambda^2\mathbb R^4\cong\mathbb R^6$ gives
$g_{ij}^2-h_{ij}=2\langle\omega_i,\omega_j\rangle^2$, so the Welch bound for $q$
unit vectors in a $6$-dimensional space yields
\[
E_4\ \ge\ \frac{q^4}{16}-\frac{q^3}{4}+\frac{q^2}{12},
\]
with equality exactly when $\{\omega_k\}$ is a tight frame of \emph{simple}
bivectors. This gives $30$, $400/3$, $1150/3$ and $876$ at $q=6,8,10,12$, matching
the values found numerically.

At $q=6$ the bound reads $E_4\ge30$, attained by the incidence family of $K_4$; and
there the tight-frame condition degenerates to ``orthonormal basis'', which is the
one value of $q$ at which it can be met by simple bivectors alone. That is why
graphs are extremal at $q=6$ and not at $q=8$, where the minimum $400/3$ lies
strictly below the graph minimum $134$: the coincidence is the degeneration, not a
principle. Even at $q=6$ the graphs are only a slice of the minimizers, the
equality locus is $4$-dimensional modulo $O(4)$, and taking the $\omega_k$ from the
six icosahedral axes, with the two $\mathbb R^3$ components Galois conjugate over
$\mathbb Q(\sqrt5)$, produces a noncommuting family with all
$g_{ij}=4/5$ and $E_4=30$. A lower bound of the same shape is known in the random setting: Brito,
Dumitriu and Harris \cite{BDH} show that a random bipartite biregular graph
satisfies $\eta^+_{\min}\ge\sqrt{d_1-1}-\sqrt{d_2-1}-\varepsilon_n$ with high
probability. That is an asymptotic statement about adjacency eigenvalues of
random graphs; Conjecture~\ref{conj:gap} is a deterministic statement about
matching polynomial roots of every graph.



\section{Mixed discriminants, the cavity recursion, and the reflection}

\subsection{The coefficients are mixed discriminants}\label{sec:mixeddisc}

Theorem~\ref{thm:matchingform} writes $M_r$ as a sum of squared Pl\"ucker norms, and
Proposition~\ref{prop:bridge} identifies the coefficients with matching numbers on the commuting
locus. Off that locus they move continuously and remain nonnegative and log-concave. They are not
a new invariant: they are mixed discriminants, the object the mixed characteristic polynomial is
assembled from.

Recall the mixed discriminant of $p$ symmetric $p\times p$ matrices,
\[
D(B_1,\dots,B_p)=\frac1{p!}\sum_{\sigma\in S_p}
\det\bigl[B_{\sigma(1)}\varepsilon_1\ \big|\ \cdots\ \big|\ B_{\sigma(p)}\varepsilon_p\bigr],
\]
introduced by Aleksandrov \cite{Alek} in the theory of mixed volumes. Writing
$\mu(x)=\sum_sc_sx^{\,p-s}$ for a tight rank-$b$ family $A_1,\dots,A_q$ on $\mathbb R^p$,
multilinearity of the determinant gives
\begin{equation}\label{eq:mixeddisc}
c_s=(-1)^s\,s!\,\binom ps\sum_{|S|=s}D\bigl(A_k:k\in S,\ I^{\,p-s}\bigr),
\end{equation}
checked to $1.4\times10^{-14}$ at every $s$ on five families with commutator norms between $0.84$
and $1.00$, so on families far from commuting (\texttt{code/mixeddisc.py}).

Identity~\eqref{eq:mixeddisc} says how to compute the coefficients, not what they are. The following
says what they are, and it needs no new object either.

\begin{proposition}[the coefficients are weighted partial transversals]\label{prop:transversal}
Let $A_1,\dots,A_q$ be rank-$b$ orthogonal projections on $\mathbb R^p$ and let
$u_{k,1},\dots,u_{k,b}$ be an orthonormal basis of $\operatorname{ran}A_k$. Writing
$\mu(y)=\sum_sm_sy^{\,p-s}$,
\begin{equation}\label{eq:transversal}
m_s=(-1)^s\sum_{|S|=s}\ \sum_{r:S\to[b]}\det\operatorname{Gram}\bigl(u_{k,r(k)}:k\in S\bigr).
\end{equation}
\end{proposition}

\begin{proof}
Expanding a determinant perturbed by rank ones,
$\det(yI+\sum_{k,r}w_{k,r}u_{k,r}u_{k,r}^{\mathsf T})
=\sum_Ty^{\,p-|T|}\bigl(\prod_{(k,r)\in T}w_{k,r}\bigr)\det\operatorname{Gram}(T)$,
the sum being over sets $T$ of pairs. Put $w_{k,r}=z_k$, so that the monomial attached to $T$ is
$\prod_kz_k^{t_k}$ with $t_k$ the number of pairs of $T$ in block $k$. Now
$(1-\partial_z)z^t$ evaluated at $z=0$ is $1$ for $t=0$, $-1$ for $t=1$ and $0$ for $t\ge2$, so
applying $\prod_k(1-\partial_{z_k})$ and setting $z=0$ kills every $T$ using two or more vectors from
one block and leaves \eqref{eq:transversal}.
\end{proof}

The inner object is a partial transversal of the block system, one representative chosen from each of
$s$ blocks, weighted by the squared volume they span. On the commuting locus the ranges are
coordinate subspaces, the representatives are standard basis vectors, and the squared volume is $1$
when they are distinct and $0$ when two coincide: the count of partial systems of distinct
representatives, which by Proposition~\ref{prop:bridge} is the matching number of the incidence graph.
Off the locus the weight is the continuous measure of how independent the representatives are. So the
object being weighted off the locus is exactly the object being counted on it. Verified against $\mu$
coefficient by coefficient up to $s=6$ at $C_4$, $K_{3,3}$ and the Fano family, commuting and
deformed (\texttt{code/transversal.py}); no novelty is claimed, the rank-one case being the standard
multi-affine picture of \cite{MSS}. The step the proof turns on, that $(1-\partial_z)z^t$ vanishes at
the origin from $t=2$ on, is machine-checked as
\texttt{Transversal.one\_sub\_deriv\_eval\_zero}, and the two trace identities of
Proposition~\ref{prop:leadrigid} as \texttt{Transversal.sum\_trace\_pairs} and
\texttt{sum\_trace\_triples}, with \texttt{trace\_three\_symm} for the fact that every ordering of a
triple of symmetric blocks carries the same trace. Neither proposition is formalised whole, Mathlib
carrying no mixed characteristic polynomial.

Expanding the Gram determinant makes the rigidity of the leading coefficients a computation rather
than an observation, and in a stronger form than Proposition~\ref{prop:rigid}, which moves along one
rotation.

\begin{proposition}[the four leading coefficients are class-wide constants]\label{prop:leadrigid}
For every tight rank-$b$ projection family at fixed $(n,q,a,b)$,
\[
m_0=1,\quad m_1=-qb,\quad m_2=\binom q2b^2-\tfrac{a^2n-qb}2,\quad
m_3=-\Bigl[\binom q3b^3-\tfrac{b(q-2)(a^2n-qb)}2+\tfrac{a^3n-3a^2n+2qb}3\Bigr].
\]
\end{proposition}

\begin{proof}
The Gram determinant of $s$ unit vectors is $1$ for $s=1$; $1-\langle u,v\rangle^2$ for $s=2$; and
$1-\langle u,v\rangle^2-\langle u,w\rangle^2-\langle v,w\rangle^2
+2\langle u,v\rangle\langle v,w\rangle\langle w,u\rangle$ for $s=3$. Summing over $r$ turns each
closed walk of inner products into a trace, so $m_2$ involves only $\sum_{k<l}\operatorname{tr}A_kA_l$
and $m_3$ only that together with $\sum_{k<l<m}\operatorname{tr}A_kA_lA_m$. Both are fixed by the
hypotheses: from $\sum_kA_k=aI$, $\operatorname{tr}A_k=b$ and $A_k^2=A_k$,
$\operatorname{tr}\bigl((\sum_kA_k)^2\bigr)=a^2n$ gives $\sum_{k<l}\operatorname{tr}A_kA_l
=(a^2n-qb)/2$, and $\operatorname{tr}\bigl((\sum_kA_k)^3\bigr)=a^3n$ gives
$\sum_{k<l<m}\operatorname{tr}A_kA_lA_m=(a^3n-3a^2n+2qb)/6$, every ordered triple of distinct indices
contributing the same value because the blocks are symmetric, so that
$\operatorname{tr}A_kA_mA_l=\operatorname{tr}A_kA_lA_m$.
\end{proof}

At $s=4$ the Gram determinant carries the four-cycle
$\langle u_1,u_2\rangle\langle u_2,u_3\rangle\langle u_3,u_4\rangle\langle u_4,u_1\rangle$, whose sum
over $r$ is $\operatorname{tr}A_kA_lA_mA_r$ with four distinct indices, and the single relation from
$\operatorname{tr}\bigl((\sum_kA_k)^4\bigr)=a^4n$ cannot pin the three independent cyclic classes. So
$m_4$ is the first coefficient carrying a genuine joint invariant, which is why the cut falls at four.
It does move: at the Fano family $m_4$ ranges over an interval of length $6.03$ across random tight
families against the commuting value $1218$, and at $\mathrm{PG}(2,3)$ the four leading coefficients
sit at $1,-52,1170,-14976$ for commuting and deformed families alike while $m_4$, $m_5$ and $m_6$
move from $120627$, $-638820$, $2258100$ (\texttt{code/transversal.py}).

The cycle expansion has a combinatorial reading, and it is one substitution.

\begin{proposition}[the coefficients count systems of distinct representatives]\label{prop:sdr}
For a coordinate family, and for every $S$,
\[
\sum_{\sigma\in\mathrm{Sym}(S)}\mathrm{sgn}(\sigma)\prod_{\text{cycles }c}\Bigl|\bigcap_{k\in c}e_k\Bigr|
=\#\bigl\{f:S\to V \ :\ f(k)\in e_k,\ f\text{ injective}\bigr\},
\]
so $|m_s|$ is the number of partial systems of distinct representatives of size $s$.
\end{proposition}

\begin{proof}
The permutations whose cycle partition is a given $\pi$ contribute
$\prod_{B\in\pi}(-1)^{|B|-1}(|B|-1)!=\mu(\hat0,\pi)$ in total, the M\"obius function of the
partition lattice. The product $\prod_{B\in\pi}|\bigcap_{k\in B}e_k|$ counts the choice functions
constant on every block of $\pi$, so the sum is the M\"obius inversion over the lattice of collisions
and returns the choice functions with no collision.
\end{proof}

This is the classical inclusion-exclusion count of systems of distinct representatives in its
partition-lattice form, and no novelty is claimed for it; it is recorded because of what happens
next. Off the commuting locus the expansion is unchanged except that
\[
\Bigl|\bigcap_{k\in B}e_k\Bigr|\qquad\text{becomes}\qquad \mathrm{tr}\Bigl(\prod_{k\in B}A_k\Bigr).
\]
So $\mu$ is the same signed count for a family of \emph{subspaces}, with the number of points common
to a set of blocks replaced by the cyclic trace of their projections. Three facts recorded separately
above have that single source. The rigidity of $m_0,\dots,m_3$ is that blocks of size at most three
have their cyclic traces fixed by $\sum_kA_k=aI$ and idempotency. The cut at four is that
$\mathrm{tr}(A_iA_jA_kA_l)$ depends on the cyclic order, of which there are three on four blocks,
where an intersection number has no order at all: at a $4$-subset the three orders give one value on
the locus and three off it, spread $0$ against $0.24$ at $K_{3,3}$
(\texttt{code/sdr.py}). And the weights being multi-way rather than pairwise is why no edge-weighted
graph carries them. Verified on $530$ subsets across four families on the locus, and against $\mu$ up
to $s=5$ off it.

The second-order deformation is carried entirely by the four-block data, and that is the structural
statement behind everything above. With $A_k=P_k+\varepsilon D_k+\varepsilon^2X_k$, every cyclic
trace variation is built from the products $D_k(u,v)D_l(u,v)$: two $D$'s give
$\mathrm{tr}(D_xBD_yC)=\sum_{u,v}B(v)C(u)D_x(u,v)D_y(u,v)$ with $B,C$ diagonal, and one $X$ enters
only through its diagonal, which order two forces. So all the variations live in one space, cut by
tightness and the cone condition, and one can ask how much of it each block size spans. The answer is
that nothing beyond four contributes: the rank of the data from blocks of size at most four equals
that from size at most five, six and seven, at $C_6$, $K_{3,3}$, the Fano family and the cube, with
ranks $22$, $82$, $57$ and $189$ (\texttt{code/relations.py}). The relations among the variations are
the cokernel of that map and number $22$, $65$, $28$ and $148$, against the two coming from $P_2$ and
$P_3$ being constant.

That settles the shape of the whole picture. Below four the cyclic traces are fixed by tightness and
idempotency, so $m_0,\dots,m_3$ cannot move. At four the intersection number acquires a cyclic order
and splits into three, which is the only new information the deformation carries. Above four nothing
further enters, so $\delta m_5$ and $\delta m_4$ are two functionals on the same space and can be
proportional, while $\delta m_6$ is a third functional on it and need not be.

Half of the ceiling is a counting argument rather than a computation.

\begin{proposition}[long cycles do not contribute]\label{prop:arcbound}
The second-order row of a cyclic word of length $m$ vanishes identically whenever $m>2a$.
\end{proposition}

\begin{proof}
A nonzero coefficient at $z^{uv}_{xy}$ requires $u\in I_1$ and $v\in I_2$, or the same with $u,v$
exchanged. Every block on the first arc then contains $u$, and $x$ and $y$ both split the pair
$\{u,v\}$, so whichever of them contains $u$ is a further block containing $u$ and not on that arc.
A point of a tight family lies in exactly $a$ blocks, so the first arc has at most $a-1$ of them, and
the second at most $a-1$ likewise; if instead $x$ and $y$ both contain $u$ the first arc has at most
$a-2$ and the second at most $a$. Either way $m=2+|\text{arc}_1|+|\text{arc}_2|\le2a$. The diagonal
terms need every other block to contain the same point and so require $m\le a+1$.
\end{proof}

The counting step is machine-checked as \texttt{ArcBound.card\_le\_of\_mem\_all} and
\texttt{arc\_bound}. The bound is tight: at $K_{3,3}$ and the cube, both with $a=3$, rows of length
exactly $6=2a$ are nonzero (\texttt{code/ceiling.py}). At $a=2$ it reads $m\le4$ and closes the
four-block ceiling outright, with no linear algebra. At $a\ge3$ it leaves rows of length five and six
nonzero, and their dependence on the shorter ones, which the ranks above exhibit, is not explained
here. That dependence is not local either: a length-five row lies in the span of the rows on its own
five blocks at $K_{3,3}$ and at the cube, and does not at the Fano family, so no single local identity
accounts for it and a proof for $a\ge3$ has to be global.

At the level of these rows the second half of Conjecture~\ref{conj:lockstep} is not conjectural. It
follows from the shape of a single difference, and the mechanism is that the difference lies in the
ideal of functionals that tightness makes identically zero.

\begin{proposition}[the five-block row is determined by the four-block row]\label{prop:I2}
Write $W_4^{\mathrm{row}}$ and $W_5^{\mathrm{row}}$ for the second-order rows of $W_4$ and $W_5$ in
the coordinates $z^{uv}_{kl}=D_k(u,v)D_l(u,v)$, and suppose
\[
W_5^{\mathrm{row}}-4(a-2)\,W_4^{\mathrm{row}}=c\,T,\qquad
T(z)=\sum_{u<v}\ \sum_{k,l\in K(u,v)}z^{uv}_{kl},
\]
for some scalar $c$. Then $\delta W_5=4(a-2)\,\delta W_4$ for every direction and every curve.
\end{proposition}

\begin{proof}
$T$ is the vanishing functional of tightness. Indeed
\[
T(z)=\sum_{u<v}\ \sum_{k,l}D_k(u,v)D_l(u,v)=\sum_{u<v}\Bigl(\sum_kD_k(u,v)\Bigr)^2,
\]
and $\sum_kD_k=0$ is the linearisation of $\sum_kA_k=aI$, so $T$ vanishes at every admissible $z$.
Hence $W_5^{\mathrm{row}}$ and $4(a-2)W_4^{\mathrm{row}}$ agree as functionals on the accessible data,
which is the only place either is evaluated.
\end{proof}

The hypothesis is a statement about coefficient patterns and is checked directly: the difference takes
exactly one value on the diagonal coordinates $z^{uv}_{kk}$ and exactly twice that on the off-diagonal
$z^{uv}_{kl}$, the doubling being the $k<l$ storage convention, at $C_6$, $C_8$, $K_{3,3}$, the Fano
family and the cube (\texttt{code/ceiling.py}). The scalar is not derived. It is $0$, $-12$ and $-60$
at $a=2,3,4$, the same at $q=7$, $9$ and $12$ and so independent of $q$; a prediction of $-12(a-2)$
frozen before the $a=4$ computation missed, giving $-24$ against $-60$. None of that matters for the
proposition, whose conclusion holds whatever $c$ is.

The first half of the conjecture, $\delta R=2(a-2)\delta Q_2$, has no such argument. It holds exactly,
to between $10^{-15}$ and $4\times10^{-14}$ at all four families, as a certificate: the functional
$\phi=(d_P-2(a-2)e_P;\,c_T)$ annihilates the image of the $z$-map. But its difference is not a
multiple of $T$. That shape holds at the Fano family and $K_{3,3}$, with $c=-6$ and $-3$ respectively,
and fails at $C_6$ and the cube, where the difference takes several values on each kind of coordinate;
the weaker form $\sum_{u<v}\alpha_{uv}T^{uv}$, one coefficient per vertex pair, fails at the same two.
So the first half must also engage the quadrics $Q_j$, and it rests on a per-family certificate.
Conjecture~\ref{conj:lockstep} is therefore stated as a conjecture, with its second half proved.

Among the coefficients that do move, the first two are not independent, and this is the sharpest
statement the class supports.

\begin{proposition}[the second-order variation of $m_5$]\label{prop:lockstep}
Let $A_0$ be a commuting tight rank-$b$ projection family and $A(\varepsilon)$ a curve in the variety
with $A(0)=A_0$. Every coefficient of $\mu$ varies to order $\varepsilon^2$, the first order
vanishing because the blocks are diagonal and the tangent directions off-diagonal, and
\[
\delta m_5=-b(q-4)\,\delta m_4+2\bigl(\delta R-\delta W_5\bigr),
\]
with $R$ the sum over disjoint triple-and-pair of $\tau\,t$ and $W_5$ the five-cycle trace sum.
\end{proposition}

\begin{conjecture}[the first two moving coefficients are locked]\label{conj:lockstep}
$\delta R-\delta W_5=2(a-2)\,\delta m_4$, so that $\delta m_5=c\,\delta m_4$ with
$c=4(a-2)-b(q-4)$ an integer independent of the direction and of the curve.
\end{conjecture}

Proposition~\ref{prop:lockstep} is the derived part. Expanding the cycle formula at $s=5$ and
summing over $5$-subsets,
\[
m_5=-\Bigl[\tbinom q5b^5-b^3\tbinom{q-2}3P_2+b(q-4)Q_2+2b^2\tbinom{q-3}2P_3-2R-2b(q-4)W_4+2W_5\Bigr],
\]
with $R$ the sum over disjoint triple-and-pair of $\tau\,t$ and $W_5$ the five-cycle trace sum. Since
$P_2$ and $P_3$ are class-wide constants and $\delta m_4=\delta Q_2-2\delta W_4$,
\[
\delta m_5=-b(q-4)\,\delta m_4+2\bigl(\delta R-\delta W_5\bigr).
\]
which is the proposition. A universal $c$ therefore requires $\delta R-\delta W_5$ to be a universal
multiple of $\delta m_4$, which is Conjecture~\ref{conj:lockstep} and is measured rather than proved.
The evidence: the ratio is direction-independent to between $10^{-9}$ and $10^{-7}$ at six
families, and unchanged when the second-order term of the curve is perturbed at scales $1$ and $3$,
so it is a property of the variety and not of the retraction. The formula was fitted on $C_6$, $C_8$,
$K_{3,3}$ and the Fano family, giving $-4$, $-8$, $-6$, $-5$, and then predicted $0$ for the
$2$-regular $3$-uniform family on six vertices and $-16$ for $\mathrm{AG}(2,3)$ before those were
computed; both were hit, at $0.00000000$ and $-15.99999999$ (\texttt{code/lockstep.py}).

The lock does not extend. The ratio $\delta m_6/\delta m_4$ varies over directions by $1.6$ at $C_6$
and $3.6$ at $K_{3,3}$, so the second-order deformation is constrained but not one-dimensional. Taken
with Proposition~\ref{prop:leadrigid}, the second-order variation of $\mu$ lies in a subspace of
codimension at least five in the $(n+1)$-dimensional coefficient space, for every family and every
direction.

One reading is excluded outright. A weighted matching polynomial would have the transversal weight
factor through pairwise data, and it does not, for a reason visible on the commuting locus with no
deformation. In the Fano plane any two lines meet in one point, so every triple of blocks carries the
same pairwise traces; but three lines are either concurrent or form a triangle, and
$\operatorname{tr}P_kP_lP_m$ is $1$ or $0$ accordingly. Two triples with identical pairwise weights
and different transversal weights is exactly the failure of pairwise factorisation, and the same
happens at $\mathrm{PG}(2,3)$. The weight is irreducibly higher than pairwise, and no edge-weighted
graph carries it.

Two properties of $D$ have to hold before \eqref{eq:mixeddisc} even parses, and both are now
machine-checked rather than carried as prose. The sum on the right is over \emph{sets} $S$, which is
legitimate only because $D$ is symmetric in its arguments
(\texttt{MixedDiscriminant.mixedDisc\_comp\_perm}); and calling $D$ the multilinear extension of the
determinant is a claim whose two halves are $D(B,\dots,B)=\det B$
(\texttt{mixedDisc\_const}) and additivity in each argument separately
(\texttt{mixedDisc\_update\_add}), the latter holding because in the term indexed by $\sigma$ the
matrix $B_k$ occupies exactly one column, the one numbered $\sigma^{-1}(k)$. Aleksandrov's
inequality, Bapat's theorem and Gurvits' theorem are cited and not formalised, and neither is
\eqref{eq:mixeddisc} itself.

Two properties recorded earlier as observations are then consequences of classical results rather
than facts about this construction. The $c_s$ alternate in sign because a mixed discriminant of
positive semidefinite arguments is nonnegative, which is Aleksandrov's inequality; the
quantitative theory of that quantity is Bapat's \cite{Bapat}, whose conjectured sharp lower bound
for doubly stochastic tuples was proved by Gurvits \cite{Gurvits}. And the coefficients are
log-concave because $\mu$ is real-rooted, by \cite{MSS}, so Newton's inequalities apply; the
measured ratios run from $1.07$ to $1.20$.

The same identification places the note's own obstruction in classical territory. A tight family
with $q=p$ and $a=b$, scaled by $1/a$, has $\sum_kA_k=I$ and $\operatorname{tr}A_k=1$, so it is a
doubly stochastic tuple in the sense of \cite{Gurvits}, and its mixed discriminant is at least
$p!/p^p$ with equality only at $A_k=I/p$. That tuple is exactly the one recorded in
\S\ref{sec:leaves} as satisfying rank, trace and $\sum_kA_k=aI$ while exceeding the band: its
greatest root is $4.698$ at $p=4$ against a band edge of $4$, and it grows with $p$. The example
that blocks every argument reading only those three data is therefore the van der Waerden extremal,
not an ad hoc construction. Both halves are checked in \texttt{code/mixeddisc.py}, in the
normalisation with $D(I,\dots,I)=p!$.

What \eqref{eq:mixeddisc} settles is the question of what the coefficients count away from the
commuting locus. On the locus, the mixed discriminant of coordinate projections is a count, of
systems of distinct representatives, which is Proposition~\ref{prop:bridge}. Off it, it is the
multilinear extension of that count and is not itself a count. No novelty is claimed for
\eqref{eq:mixeddisc}; it is recorded because the alternative was to keep describing a classical
object in ad hoc terms.


\section{The reflection, and why the band edges are never roots}

One structural fact explains why $b=2$ is the case that closes, and it is a
reflection. The subdivision identity
$\mu_{S(H)}(x)=x^{|E|-|V|}\,\mathcal{LM}(H,x^2)$ is due to Yan and Yeh
\cite{YY}, who also treat the regular and semiregular cases; see
\cite{WWM}, Corollary 2.3, for a restatement through the multivariate matching
polynomial. For $H$ $D$-regular it reads $f_G(y)=\mu_H(y-D)$, and
$\mu_H(-z)=(-1)^{|V(H)|}\mu_H(z)$, so
\begin{equation}\label{eq:reflect}
f_G(2c-y)=(-1)^{p}f_G(y),\qquad c=(a-1)+(b-1),
\end{equation}
and $c$ is the midpoint of $[(s-t)^2,(s+t)^2]$ because $(s\pm t)^2=c\pm2\sqrt m$.
The roots of $f_G$ are therefore symmetric about the centre of the band, so for
$b=2$ the lower bound in \eqref{eq:twosided} is \emph{equivalent} to the upper
one, and Theorem~\ref{thm:subdiv} is that equivalence. The reflection
\eqref{eq:reflect} fails as soon as both degrees exceed $2$: it fails for
$K_{3,4}$, $K_{3,5}$ and already for the regular $K_{3,3}$, so regularity does
not rescue it. Since $b=2$ holds exactly when $G$ is the subdivision of an
$a$-regular multigraph, \eqref{eq:reflect} describes the whole of the
phenomenon, and for $b\ge3$ the lower bound is not a repackaging of the upper
bound but carries information the upper bound does not. We note that
\eqref{eq:reflect} has no adjacency analogue: the companion identity
$\varphi(S(H),x)=x^{|E|-|V|}\varphi(H,x^2-D)$ holds as well, but $\varphi(H,\cdot)$
has a parity only for bipartite $H$, whereas $\mu_H$ always does. For the
Petersen graph the matching-side values are symmetric about $D=3$ and the
adjacency-side values are not.

A consequence of the first is that the band edges are never roots, by an
argument of a different kind.

\begin{proposition}\label{prop:galois}
Let $G$ be $(a,b)$-biregular bipartite with $a,b\ge2$, and suppose
$m=(a-1)(b-1)$ is not a perfect square. Then neither $(s-t)^2$ nor $(s+t)^2$ is
a root of $f_G$.
\end{proposition}

\begin{proof}
$f_G$ has integer coefficients and $(s\mp t)^2=(a+b-2)\mp2\sqrt m$, so for $m$
not a square the two band edges are conjugate over $\mathbb Q$; if one is a root
of $f_G$, so is the other. Heilmann and Lieb already give a \emph{strict} bound,
$|a|<2\sqrt{B_1}$ with $B_1=\Delta-1$ for unit edge weights \cite[Thm.~4.3]{HL},
but at the max-degree radius, which for $a\ne b$ exceeds $s+t$; strictness at the
biregular radius therefore has to be argued separately. Every root of $\mu_G$ is
bounded by the spectral
radius of a path tree $T_u(G)$ \cite{Go}, and $T_u(G)$ is finite and embeds
in the universal cover of $G$, the $(a,b)$-biregular tree $T_{a,b}$, with
degrees bounded by those of $T_{a,b}$. Now $\rho(T_{a,b})=s+t$ is the supremum
of $\rho$ over finite subtrees of $T_{a,b}$, and that supremum is not attained:
any finite subtree is a proper subgraph of a larger connected finite subtree,
and $\rho$ is strictly monotone under proper subgraph containment for finite
connected graphs. Hence $\rho(T_u(G))<s+t$, so $(s+t)^2$ is not a root of
$f_G$, and therefore neither is $(s-t)^2$.
\end{proof}

The hypothesis on $m$ is needed for the argument, not merely for convenience:
when $m$ is a square the two band edges are rational and independent, and
conjugacy says nothing. The smallest such case is $(a,b)=(5,2)$, where $m=4$
and the edges are $1$ and $9$. There the conclusion nevertheless holds, but for
a different reason, the Heilmann--Lieb bound is strict for finite graphs, so
Theorem~\ref{thm:subdiv} already excludes the lower edge. Proposition~\ref{prop:galois}
is the statement for the cases those methods do not reach.

Evaluating $f_G$ in $\mathbb Z[\sqrt m]$ confirms this and exhibits the
conjugacy: for $K_{3,4}$, $K_{3,5}$, $K_{3,6}$, $K_{4,5}$ the two edge values
are $53\mp30\sqrt6$, $72\mp40\sqrt8$, $91\mp50\sqrt{10}$, $1877\mp512\sqrt{12}$.
For $b=2$ the irrational parts vanish, as \eqref{eq:reflect} requires.
Proposition~\ref{prop:galois} is worth isolating because it uses no induction on
vertex deletion, and so is not subject to the obstruction of
Section~\ref{sec:must}; it is
the only tool we have found with that property, although it controls only the
edges and not the interior of the gap.


\section{The multivariate polynomial and the cavity recursion}

The polynomials $g$ admit a clean description. Since a $k$-matching leaves $p-k$
vertices of $A$ unmatched,
\begin{equation}\label{eq:multivar}
f_G(y)=\sum_{k}(-1)^k m_k\,y^{p-k}
=\sum_{M}(-1)^{|M|}\prod_{v\in A\setminus V(M)}y
=\mathcal M\bigl(G;\,x_A=y,\ x_B=1\bigr),
\end{equation}
the multivariate matching polynomial of Heilmann and Lieb under the asymmetric
specialization that puts the variable on one side of the bipartition and $1$ on
the other. So the $g$ are not exotic: they are shifts of matching polynomials in
the multivariate sense. This is a reformulation and not an advance: for
bipartite $G$ every edge of a matching covers one vertex of each side, so
$\mathcal M(G;x_A=y,x_B=z)$ depends only on the product $yz$, and the asymmetric
specialization therefore carries exactly the information of the univariate
matching polynomial. In particular no quantitative Heilmann--Lieb theorem for
non-uniform specializations can help here, and none is available: the
multivariate bounds of \cite{HL} apply a single graph-wide threshold to every
variable.

Deleting a vertex now gives two different recursions, because only $A$-vertices
carry the variable:
\[
v\in A:\quad f_G=y\,f_{G-v}-\!\!\sum_{u\sim v}\!f_{G-v-u},
\qquad
u\in B:\quad f_G=f_{G-u}-\!\!\sum_{v\sim u}\!f_{G-u-v},
\]
whence the cavity recursion
$R_A=y-\sum1/R_B$, $R_B=1-\sum1/R_A$. On the $(a,b)$-biregular tree, with
$P=a-1$, $Q=b-1$, its fixed point satisfies $R^2+(P-Q-y)R+yQ=0$, of
discriminant
\[
D(y)=(P-Q-y)^2-4yQ,
\qquad
D(y)=0\iff y=(\sqrt P\pm\sqrt Q)^2 .
\]
The two roots are exactly the spectrum edges in the variable $y=x^2$, and $D<0$
strictly between them: the fixed point is complex precisely on the spectrum and
real precisely off it, as it must be. Both recursions and the discriminant have
been checked exactly.

Conjecture~\ref{conj:gap} for biregular $G$ is therefore the statement that the
asymmetric recursion, started from the leaves of the path tree, never produces
$R=0$ while $D>0$. The natural invariant region is
\[
R_A\le-\alpha<0,\qquad R_B\in[\,1,\ 1+Q/\alpha\,],
\]
for a parameter $\alpha>0$. A $B$-vertex with children in the region has
$R_B=1+\sum1/|R_A|\in(1,1+Q/\alpha]$, and a $B$-leaf carries $R_B=1$, exactly
the lower boundary. An $A$-vertex with its full complement of $P$ children has
$R_A\le y-P\alpha/(\alpha+Q)$, so the region is closed precisely when
$y\le h(\alpha)$, where $h(\alpha)=\alpha(P-Q-\alpha)/(\alpha+Q)$. The
following is elementary and is machine-checked in Lean~4
(\texttt{region\_closure\_bound}, \texttt{gap\_edge\_identity}):
\begin{equation}\label{eq:square}
(s-t)^2(\alpha+t^2)-\alpha(s^2-t^2-\alpha)=\bigl(\alpha-t(s-t)\bigr)^2 ,
\end{equation}
so $h(\alpha)\le(s-t)^2$ for every $\alpha$, with equality at
$\alpha=t(s-t)=\sqrt{PQ}-Q$. The region therefore closes for every $y$ strictly
below the gap edge and for no $y$ above it: the invariant region is calibrated
exactly to the spectrum, and \eqref{eq:square} is the reason.


\section{The formalization}

Everything asserted in this paper is machine-checked in Lean~4/Mathlib unless labelled
otherwise. The sources are \texttt{RamaLean/} in the \texttt{rama-notes} repository; every file
carries its own statement of what is proved and what is assumed, and the standing blocker
throughout is that Mathlib has no matching polynomial and no mixed characteristic polynomial, so
the identities connecting the combinatorial and operator forms are stated rather than derived.
The full file-by-file record, with the axiom base of each theorem, is in \cite{Supp2}.



\begin{thebibliography}{9}
\bibitem{Alek} A.~D.~Aleksandrov, \emph{On the theory of mixed volumes of convex bodies IV.
Mixed discriminants and mixed volumes}, Mat.\ Sb.\ (N.S.) \textbf{3} (1938), 227--251.
\bibitem{Bapat} R.~B.~Bapat, \emph{Mixed discriminants of positive semidefinite matrices},
Linear Algebra Appl.\ \textbf{126} (1989), 107--124.
\bibitem{Gurvits} L.~Gurvits, \emph{The van der Waerden conjecture for mixed discriminants},
Adv.\ Math.\ \textbf{200} (2006), 435--454.
\bibitem{Zhang} J.-C.~Wan, Y.~Wang, Y.-Z.~Fan, \emph{A hypergraph Heilmann--Lieb
theorem}, arXiv:2206.09558 (2022).
\bibitem{XXZ} Z.~Xu, Z.~Xu, Z.~Zhu, \emph{Improved bounds in Weaver's
$\mathrm{KS}_r$ conjecture for high rank positive semidefinite matrices},
arXiv:2106.09205 (2021).
\bibitem{GG} C.~Godsil, I.~Gutman, \emph{On the matching polynomial of a
graph}, in: Algebraic Methods in Graph Theory, 1981.
\bibitem{HPS} C.~Hall, D.~Puder, W.~Sawin, \emph{Ramanujan coverings of
graphs}, Adv.\ Math.\ \textbf{323} (2018), 367--410.
\bibitem{CGHRS} G.~Cochran, C.~Groothuis, A.~Herring, R.~Rohatgi,
E.~Stucky, \emph{A new (combinatorial) proof of the commutativity of
matching polynomials for cycles}, arXiv:1810.05889 (2018). Preprint; no journal
reference or DOI as of August 2026.
\bibitem{WWM} J.~Wan, W.~Wang, A.~Mohammadian, \emph{On the matching
polynomials of graphs with small number of cycles}, arXiv:2103.10580 (2021).
\bibitem{SFM} Y.-M.~Song, Y.-Z.~Fan, Z.~Miao, \emph{Hypergraph coverings and
Ramanujan hypergraphs}, arXiv:2310.01771 (2023).
\bibitem{YY} W.~Yan, Y.-N.~Yeh, \emph{On the matching polynomial of subdivision
graphs}, Discrete Appl.\ Math.\ \textbf{157} (2009) 195--200. The identity used
here is their Corollary 2.1: for $G$ $d$-regular, $m(S(G),x)=x^{m-n}m(G,x^2-d)$.
\bibitem{BDH} G.~Brito, I.~Dumitriu, K.~D.~Harris, \emph{Spectral gap in random
bipartite biregular graphs and applications}, Combin.\ Probab.\ Comput.\
\textbf{31} (2022) 229--267.
\bibitem{HL} O.~J.~Heilmann, E.~H.~Lieb, \emph{Theory of monomer-dimer
systems}, Comm.\ Math.\ Phys.\ \textbf{25} (1972) 190--232.
\bibitem{CC} V.~Y. Chernyak, M.~Chertkov, \emph{Fermions and loops on graphs.
II. A monomer-dimer model as a series of determinants}, J. Stat. Mech. (2008)
P12012.

\bibitem{WC} Y.~Watanabe, M.~Chertkov, \emph{Belief propagation and loop calculus
for the permanent of a non-negative matrix}, J. Phys. A \textbf{43} (2010)
242002.

\bibitem{Je} M.~Jerrum, \emph{Two-dimensional monomer-dimer systems are
computationally intractable}, J. Statist. Phys. \textbf{48} (1987) 121--134.

\bibitem{WF} Y.~Watanabe, K.~Fukumizu, \emph{New graph polynomials from the Bethe
approximation of the Ising partition function}, Combin. Probab. Comput.
\textbf{20} (2011) 299--320.

\bibitem{BC} C.~Bordenave, B.~Collins, \emph{Eigenvalues of random lifts and
polynomials of random permutation matrices}, Ann.\ of Math.\ \textbf{190}
(2019) 811--875.
\bibitem{MF} J.~P.~McSorley, P.~Feinsilver, \emph{Multivariate matching
polynomials of cyclically labelled graphs}, Discrete Math.\ \textbf{309} (2009)
3205--3218.
\bibitem{MS} H.~Mizuno, I.~Sato, \emph{Characteristic polynomials of some graph
coverings}, Discrete Math.\ \textbf{142} (1995) 295--298.
\bibitem{DFMRS} C.~Dalf\'o, M.~A.~Fiol, M.~Miller, J.~Ryan, J.~\v{S}ir\'a\v{n},
\emph{An algebraic approach to lifts of digraphs}, Discrete Appl.\ Math.\
\textbf{269} (2019) 68--76.
\bibitem{MSSII} A.~W. Marcus, D.~A. Spielman, N.~Srivastava,
\emph{Interlacing families II: mixed characteristic polynomials and the Kadison--Singer
problem}, Ann. of Math. (2) \textbf{182} (2015), 327--350.

\bibitem{MSS} A.~Marcus, D.~Spielman, N.~Srivastava, \emph{Interlacing
families I: Bipartite Ramanujan graphs of all degrees}, Ann.\ of Math.\
\textbf{182} (2015), 307--325.
\bibitem{CsF} P.~Csikv\'ari, P.~E.~Frenkel, \emph{Benjamini--Schramm continuity
of root moments of graph polynomials}, European J.\ Combin.\ \textbf{52} (2016),
part B, 302--320.
\bibitem{AH} M.~Ab\'ert, T.~Hubai, \emph{Benjamini--Schramm convergence and the
distribution of chromatic roots for sparse graphs}, Combinatorica \textbf{35}
(2015), 127--151.
\bibitem{Go} C.~D.~Godsil, \emph{Matchings and walks in graphs}, J.\ Graph
Theory \textbf{5} (1981), 285--297.
\bibitem{HallCE} C.~Hall, \emph{A simple bipartite counterexample to Conjecture 10:
  spectral localization of the ordinary matching polynomial}, verification note, personal
  communication, August 2026.
\bibitem{Avakov} E.~R.~Avakov, \emph{Extremum conditions for smooth problems with
equality-type constraints}, USSR Comput.\ Math.\ Math.\ Phys.\ \textbf{25} (1985), 24--32.

\bibitem{Arutyunov} A.~V.~Arutyunov, \emph{Optimality Conditions: Abnormal and Degenerate
Problems}, Kluwer, 2000.

\bibitem{RL} M.~Ravichandran, J.~Leake, \emph{Mixed determinants and the Kadison--Singer
  problem}, Math. Ann. \textbf{377} (2020), 511--541.
\bibitem{AFH} O.~Angel, J.~Friedman, S.~Hoory, \emph{The non-backtracking spectrum of the
  universal cover of a graph}, Trans.\ Amer.\ Math.\ Soc.\ \textbf{367} (2015), 4287--4318;
  arXiv:0712.0192.
\bibitem{BCs} F.~Bencs, P.~Csikv\'ari, \emph{Evaluations of Tutte polynomials of
regular graphs}, arXiv:2105.06798 (2021).
\bibitem{PartII} V.~Bonfioli, \emph{The biregular case of a spectral-inclusion conjecture:
  the weighted plane class and its obstructions}, 2026.
\bibitem{PartI} V.~Bonfioli, \emph{Roots of $d$-matching polynomials and the spectrum of
  the universal cover: a closed form at first Betti number one, and a refuted conjecture},
  companion paper, 2026.
\end{thebibliography}

\end{document}
